\documentclass[11pt]{article}
\usepackage[utf8]{inputenc}
\usepackage[T1]{fontenc}
\usepackage{lmodern}
\usepackage[margin=1in]{geometry}
\usepackage{graphicx}
\usepackage{booktabs}
\usepackage{longtable}
\usepackage{caption}
\usepackage{float}
\usepackage{array}
\usepackage{tabularx}
\usepackage{makecell}
\usepackage{microtype}
\usepackage[hidelinks]{hyperref}

\renewcommand{\thetable}{S\arabic{table}}
\renewcommand{\thefigure}{S\arabic{figure}}

\title{\textbf{SI Appendix: Cross-Domain Reactivation of Collective Attention in Music and Film}}
\author{Cristian Candia, Camila Utreras-Cifuentes, and Francisca Droguett\\[4pt]\small Correspondence: crcandiav@gmail.com}
\date{}

\begin{document}
\maketitle

\section*{Overview}

This SI Appendix documents the construction of the analytical data, platform coverage and missingness, the expected-attention baseline, model specifications, and the robustness and inference checks supporting the observational analyses of cross-domain reactivation. The primary unit of analysis is the Billboard song-artist record: one record per distinct song-title and performer string observed on the weekly Hot 100 chart. Spotify analyses use repeated snapshots of current platform popularity. Last.fm analyses use a single snapshot of cumulative listener reach. The evidence includes descriptive attention differences, matched observational comparisons, and temporal changes around film placement. Film placement is not assumed to be randomized or exogenous.

\section*{Data Construction and Platform Coverage}

The analytical universe contains 32,124 Billboard Hot 100 song-artist records, rebuilt directly from the weekly Hot 100 chart file: for each distinct song-title and performer string, the chart debut date, the total weeks on the Hot 100 (the longest observed chart run), and the best weekly peak position are aggregated from the weekly entries. Song titles and performer strings are preserved exactly as charted, so distinct records that would collide under aggressive punctuation-free normalization (for example, holiday songs recharting under near-identical titles, or artist spellings differing only in punctuation) remain distinct.

Film appearances were identified using soundtrack data collected from Soundtrack Collector and IMDb. During the original collection, Soundtrack Collector film titles were standardized to lower case with special characters removed to search for soundtrack listings in IMDb. The analysis uses the preserved IMDb song-film records. Billboard and soundtrack records are linked by exact matches on separately normalized song and artist keys, identifying 4,683 Billboard records with at least one matched appearance with a recorded film year. For this linkage, normalization converts strings to lower case, removes punctuation and specified collaboration connectors, and collapses whitespace. Movie-metadata and snapshot-specific exposure joins instead use exact alphanumeric keys with spaces and punctuation removed. Records without an observed match are classified as having no observed film appearance. Nonmatches can reflect incomplete soundtrack coverage, undocumented uses, or limits of exact string linkage.

In July 2026, one reviewer (Camila Utreras-Cifuentes) manually assessed a random sample of 200 unique song-artist-film links, representing 191 distinct Billboard song-artist records, drawn without replacement from the 11,446 links in the final analytical linkage. Each assessment considered the song, credited performer, and film, with recorded sources and supporting screenshots. The initial review classified 180 links as confirmed, 4 as ambiguous, and 16 as incorrect. Four flagged classifications were subsequently reconciled against the supporting soundtrack credits, distinguishing songwriter from performer identity. After reconciliation, 180 links were confirmed (90\%), 5 were ambiguous (2.5\%), and 15 were incorrect (7.5\%). The original classifications and the four changes are retained in the validation records. Observed sample precision was 90\%. This validation assesses the quality of retained links and does not estimate recall or false negatives among unlinked records. Linkage errors may affect film identity, the number of appearances, and associated box-office measures.

A separate evidence annotation distinguishes performer discrepancies from other uncertainties without changing the validation labels or analytical data. Among the 15 links classified as incorrect, 14 credit a different performer while identifying the linked artist in the composition credits. The remaining case has an expanded performer credit that includes the linked artist, which does not by itself establish a different version or a false link. The five ambiguous cases comprise two with another credited performer, one with an expanded credit, one unresolved solo-artist versus band attribution, and one with a songwriter credit but no identified performer. These distinctions concern composition and performer attribution. They do not establish exact recording identity or a separate composition-level precision estimate.

Table~\ref{tab:si-platform-coverage} reports linkage support for every analyzed platform snapshot: the four high-coverage Spotify current-popularity snapshots and the single Last.fm listener snapshot.

\begin{table}[!htbp]
\centering
\footnotesize
\caption{Platform and snapshot coverage.}
\label{tab:si-platform-coverage}
\renewcommand{\arraystretch}{1.05}
\begin{tabular}{>{\raggedright\arraybackslash}p{0.08\linewidth}>{\raggedright\arraybackslash}p{0.13\linewidth}>{\centering\arraybackslash}p{0.078\linewidth}>{\centering\arraybackslash}p{0.08\linewidth}>{\raggedright\arraybackslash}p{0.41\linewidth}}
\toprule
Platform & Snapshot & Records & Film-linked & Note \\
\midrule
Spotify & October 2016 & 17,347 & 3,825 & high-coverage snapshot \\
Spotify & July 2017 & 18,189 & 4,133 & high-coverage snapshot \\
Spotify & August 2022 & 22,105 & 4,373 & high-coverage snapshot \\
Spotify & August 2025 & 26,144 & 4,616 & high-coverage snapshot \\
Last.fm & July 2017 & 17,988 & 4,090 & single available snapshot; cumulative listener reach \\
\bottomrule
\end{tabular}
\begin{flushleft}\footnotesize
Records are Billboard song-artist records with a nonmissing outcome in each snapshot after exact normalized song-artist linkage. The Spotify analyses use the four high-coverage snapshots; Last.fm provides the single July 2017 listener snapshot.
\end{flushleft}
\end{table}


Soundtrack records were linked to movie-level data through shared movie identifiers, without additional movie-title matching. Worldwide box-office revenue and production budget were collected from Box Office Mojo during construction of the film dataset. Monetary strings are converted to numeric amounts as recorded in the preserved input, with no additional inflation or present-value adjustment in the current analysis. The continuous film-visibility measure is $\ln(1+B_{is})$, where $B_{is}$ is the maximum recorded worldwide box office among films linked to song $i$ whose release year is no later than the year of snapshot $s$. Release timing is available at annual resolution, and the preserved box-office totals do not measure revenue accrued by each snapshot date. Table~\ref{tab:si-movie-metadata-coverage} reports coverage using the common denominator of 4,683 film-linked Billboard records. Analyses using box office, budget, or genre run on the corresponding supported subsets. IMDb votes, ratings, production country, language, and scene-level soundtrack prominence are not available and are not imputed.

\begin{table}[!htbp]
\centering
\small
\caption{Movie-metadata coverage and missingness among film-linked records.}
\label{tab:si-movie-metadata-coverage}
\renewcommand{\arraystretch}{1.05}
\begin{tabular}{>{\raggedright\arraybackslash}p{0.34\linewidth}>{\centering\arraybackslash}p{0.095\linewidth}>{\centering\arraybackslash}p{0.085\linewidth}>{\centering\arraybackslash}p{0.115\linewidth}>{\centering\arraybackslash}p{0.09\linewidth}}
\toprule
Variable & Available & Missing & \makecell[b]{Film-linked\\ records} & \makecell[b]{Share\\ missing} \\
\midrule
movie metadata via exact song-artist / internal movie ID & 4,682 & 1 & 4,683 & 0.000 \\
worldwide box office & 4,272 & 411 & 4,683 & 0.088 \\
production budget & 2,581 & 2,102 & 4,683 & 0.449 \\
genre & 4,682 & 1 & 4,683 & 0.000 \\
IMDb votes / rating / country / language & 0 & 4,683 & 4,683 & 1.000 \\
\bottomrule
\end{tabular}
\begin{flushleft}\footnotesize
Movie attributes are linked through exact catalog keys and exact internal movie IDs; no fuzzy title matching is used and no values are imputed. All rows share the film-linked denominator. Analyses using box office, budget, or genre run on the corresponding supported subsets.
\end{flushleft}
\end{table}


\section*{Platform Coverage and Attention Differences}

Table~\ref{tab:si-attention-premiums} reports adjusted differences in Spotify popularity and cumulative Last.fm listener reach between film-linked records and records with no observed film appearance. Estimates are reported on the analyzed outcome scales and standardized within each platform/snapshot. Models adjust for song age, squared age, Billboard weeks, and peak chart position. The pooled Spotify model adds snapshot fixed effects, and its cluster-by-song variant appears in Table~\ref{tab:si-cluster-inference}.

\begin{table}[!htbp]
\centering
\footnotesize
\caption{Film-linked attention differences in raw and standardized units.}
\label{tab:si-attention-premiums}
\renewcommand{\arraystretch}{1.05}
\begin{tabular}{>{\raggedright\arraybackslash}p{0.08\linewidth}>{\raggedright\arraybackslash}p{0.115\linewidth}>{\raggedright\arraybackslash}p{0.12\linewidth}>{\centering\arraybackslash}p{0.075\linewidth}>{\centering\arraybackslash}p{0.10\linewidth}>{\centering\arraybackslash}p{0.075\linewidth}>{\centering\arraybackslash}p{0.095\linewidth}>{\centering\arraybackslash}p{0.065\linewidth}}
\toprule
Platform & Snapshot & Raw outcome & \makecell[b]{Raw\\ estimate} & \makecell[b]{Raw\\ 95\% CI} & \makecell[b]{SD\\ estimate} & \makecell[b]{SD\\ 95\% CI} & N \\
\midrule
Spotify & October 2016 & Spotify popularity & 12.28 & [11.78, 12.77] & 0.63 & [0.61, 0.66] & 17,347 \\
Spotify & July 2017 & Spotify popularity & 13.36 & [12.86, 13.86] & 0.66 & [0.64, 0.69] & 18,189 \\
Spotify & August 2022 & Spotify popularity & 14.29 & [13.79, 14.80] & 0.66 & [0.64, 0.69] & 22,105 \\
Spotify & August 2025 & Spotify popularity & 15.09 & [14.60, 15.58] & 0.66 & [0.64, 0.69] & 26,144 \\
Last.fm & July 2017 & log1p Last.fm listeners & 1.78 & [1.72, 1.84] & 0.74 & [0.71, 0.77] & 17,988 \\
Spotify & pooled (snapshot FE) & Spotify popularity & 13.79 & [13.54, 14.04] & -- & -- & 83,785 \\
\bottomrule
\end{tabular}
\begin{flushleft}\footnotesize
Models adjust for song age, squared age, Billboard weeks, and peak chart position; intervals use HC1 robust standard errors. Standardized (SD) estimates divide the outcome by its within-platform/snapshot standard deviation for the cross-platform comparison in the main attention-differences figure. Spotify measures current popularity and Last.fm measures cumulative listener reach. The pooled Spotify model adds snapshot fixed effects; its cluster-by-song variant appears in the SI inference table.
\end{flushleft}
\end{table}


Last.fm provides one snapshot of the raw \texttt{Listeners} field, treated as cumulative listener reach and analyzed as log1p cumulative listeners. Figure~\ref{fig:si-lastfm-age-profile} shows that mean cumulative reach among records with no observed film appearance rises across the youngest age groups before declining across older groups. This cross-sectional profile is consistent with cumulative counts and cohort composition. Last.fm provides a complementary measure of listener reach and is not used to estimate the Spotify expected-attention baseline. Because only one Last.fm snapshot exists, temporal and prior-attention matched Last.fm analyses are not estimated.

\begin{figure}[H]
\centering
\includegraphics[width=.8\linewidth]{Figures/SI/figure_si_lastfm_age_profile.pdf}
\caption{Last.fm cumulative listener reach by song age, July 2017 snapshot. The unit is the Billboard song-artist record. Points are integer-age means of log1p cumulative listeners, sized by cell support. The young-age profile is not monotonically declining. This cumulative measure is not used to estimate the Spotify expected-attention baseline.}
\label{fig:si-lastfm-age-profile}
\end{figure}

\section*{Expected-Attention Baselines and MemoryDecay Diagnostics}

The expected-attention baseline is constructed only from Spotify records with no observed film appearance, separately within each high-coverage snapshot. Songs are grouped by integer age and the canonical collective-memory decay function is fitted to the mean Spotify popularity of each age cell. Each fitted curve summarizes the cross-sectional relationship between age and attention within a snapshot. It does not estimate the longitudinal decay path of an individual song. Spotify popularity is the platform 0--100 compressed index and is not pre-logged. The logarithm enters through the fitting objective $\log(\mathrm{observed}) \sim \log S(t)$. Early-age cells receive inverse-age weights (capped at their 99th percentile) to give greater weight to the early portion of the age profile. Positivity of $N$, $p$, and $r$ and $0<q<p+r$ are enforced through a log/logit reparametrization. Every fit is restarted from the full MemoryDecay grid of starting values, keeping the best converged solution. Records with no observed film appearance provide a descriptive reference for age-specific attention, although they may experience other cultural shocks.

\begin{table}[!htbp]
\centering
\small
\caption{Never-film baseline support for the snapshot-specific MemoryDecay fits.}
\label{tab:si-memorydecay-support}
\begin{tabular}{llllll}
\toprule
Snapshot & \makecell[b]{Never-film\\ song-snapshots} & \makecell[b]{Positive\\ age cells} & \makecell[b]{Age\\ range} & \makecell[b]{Median\\ cell support} & \makecell[b]{Min\\ cell support} \\
\midrule
October 2016 & 13,522 & 59 & 0--58 & 220 & 116 \\
July 2017 & 14,056 & 60 & 0--59 & 215 & 139 \\
August 2022 & 17,732 & 65 & 0--64 & 249 & 111 \\
August 2025 & 21,528 & 68 & 0--67 & 278 & 167 \\
\bottomrule
\end{tabular}
\begin{flushleft}\footnotesize
Each curve summarizes a cross-sectional age-attention relationship within one snapshot. Songs are grouped by integer age within each snapshot, and the canonical curve is fitted to the mean Spotify popularity of each positive age cell with early-age weights $1/(t+10^{-3})$ capped at their 99th percentile.
\end{flushleft}
\end{table}


Table~\ref{tab:si-memorydecay-fits} reports the fitted snapshot-specific parameters and comparisons with exponential and log-normal alternatives fitted to the same age means with the same weights. The log-normal comparison is the only model using $\log(\mathrm{age})$. The canonical model has the lowest log-space RMSE in every snapshot. In the two more recent snapshots, the fitted decay rates nearly coincide ($p+r-q$ approaches zero). These fits are flagged in the table and included in the inversion diagnostics below. Figure~\ref{fig:si-memorydecay-model-comparison} shows the fitted curves over the age means for songs with no observed film appearance.

\begin{table}[!htbp]
\centering
\footnotesize
\caption{Snapshot-specific canonical MemoryDecay parameters and model comparison.}
\label{tab:si-memorydecay-fits}
\renewcommand{\arraystretch}{1.05}
\begin{tabular}{lllllllll}
\toprule
Snapshot & $N$ & $p$ & $r$ & $q$ & \makecell[b]{Near-coincident\\ rates} & \makecell[b]{RMSE$_{\log}$\\ canonical} & \makecell[b]{RMSE$_{\log}$\\ exponential} & \makecell[b]{RMSE$_{\log}$\\ log-normal} \\
\midrule
October 2016 & 85.20 & 0.2743 & 0.5354 & 0.0249 & no & 0.0515 & 0.0664 & 0.1259 \\
July 2017 & 84.88 & 0.2123 & 0.5232 & 0.0235 & no & 0.0659 & 0.0731 & 0.1517 \\
August 2022 & 66.38 & 0.0136 & 0.0185 & 0.0321 & yes & 0.0593 & 0.0770 & 0.1441 \\
August 2025 & 69.40 & 0.0119 & 0.0221 & 0.0340 & yes & 0.0868 & 0.1133 & 0.1954 \\
\bottomrule
\end{tabular}
\begin{flushleft}\footnotesize
Canonical fits use the log-response objective, the MemoryDecay starting grid, and early-age weighting on never-film age means. The near-coincident-rates flag identifies fits with $p+r-q<0.005$; decay-equivalent diagnostics include these curves. Exponential and log-normal comparison models are fitted to the same age means with the same weights; the log-normal comparison is the only model using log(age).
\end{flushleft}
\end{table}


\clearpage
\begin{figure}[H]
\centering
\includegraphics[width=\linewidth]{Figures/SI/figure_si_memorydecay_model_comparison.pdf}
\caption{Snapshot-specific canonical, exponential, and log-normal fits on Spotify age means for songs with no observed film appearance. Points are integer-age cell means sized by cell support. All models are fitted in log-response space with early-age weighting. The canonical biexponential provides the expected-attention baseline in every snapshot. The curves describe cross-sectional age profiles.}
\label{fig:si-memorydecay-model-comparison}
\end{figure}

For every Spotify song-snapshot with nonnegative age, excess attention is observed popularity minus the fitted snapshot-specific expected value at the song's age. A positive excess-attention level alone does not establish renewed attention. The temporal comparisons evaluate changes relative to matched controls. Decay-equivalent years translate an excess-attention contrast $\Delta$ onto the fitted curve by solving $S(t^*) = S(t) + \Delta$ for $t^* < t$ over the relevant empirical age distribution. Table~\ref{tab:si-decay-equivalent-audit} reports the underlying contrast, age support, dispersion, and share of non-estimable translations for every decay-equivalent value in the paper. Non-estimable cases occur when the shifted target exceeds the fitted young-age maximum. These translations express the same observational contrasts on another scale. They provide neither independent evidence of reactivation nor literal measures of memory recovered.

\begin{table}[!htbp]
\centering
\footnotesize
\caption{Decay-equivalent-years support and inversion diagnostics.}
\label{tab:si-decay-equivalent-audit}
\renewcommand{\arraystretch}{1.05}
\begin{tabular}{>{\raggedright\arraybackslash}p{0.30\linewidth}>{\centering\arraybackslash}p{0.078\linewidth}>{\centering\arraybackslash}p{0.065\linewidth}>{\centering\arraybackslash}p{0.068\linewidth}>{\centering\arraybackslash}p{0.12\linewidth}>{\centering\arraybackslash}p{0.082\linewidth}}
\toprule
Evidence layer & \makecell[b]{$\Delta$\\ (excess\\ points)} & N & \makecell[b]{Median\\ years} & \makecell[b]{IQR /\\ 95\% CI} & \makecell[b]{Non-\\estimable} \\
\midrule
Prior-attention matched excess premium & 2.52 & 12,095 & 3.78 & [3.01, 4.64] & 0.02\% \\
Strict temporal event contrast & 2.45 & 93 & 3.41 & [2.20, 4.94] & 0.00\% \\
Strict temporal post contrast & 2.98 & 93 & 4.11 & [2.66, 5.94] & 0.00\% \\
Lower pre-film excess-attention contrast (integrated) & 5.04 & 119 & 6.88 & [3.89, 10.22] & 0.01\% \\
\bottomrule
\end{tabular}
\begin{flushleft}\footnotesize
For each contrast $\Delta$ and song age $t$, the translation solves $S(t^*) = S(t) + \Delta$ on the fitted snapshot-specific curve and reports $t - t^*$ over the relevant empirical age distribution. Non-estimable cases occur when the shifted target exceeds the fitted young-age maximum. The lower-tercile excess-attention row reports the integrated nested-bootstrap median and 95\% CI (bootstrap-specific refitted curves) with the bootstrap-wide inversion failure rate.
\end{flushleft}
\end{table}


\section*{Matched Song-Level Comparisons}

Matched song-level comparisons compare film-linked songs with songs with no observed film appearance and similar chart history, age, artist-level Billboard success, and, where available, prior Spotify attention. Matching is nearest-neighbor with replacement on standardized Euclidean distance. Prior-attention designs additionally impose a five-point caliper on popularity in the previous high-coverage snapshot. This prior observation need not precede the first film appearance. Table~\ref{tab:si-song-matching-estimates} reports the matched support and estimates behind the main matching panels. Table~\ref{tab:si-matching-balance} consolidates covariate-level balance across the song-level, strict-temporal, and prior-attention matched designs.

\begin{table}[!htbp]
\centering
\footnotesize
\caption{Matched song-level attention premiums by snapshot and design.}
\label{tab:si-song-matching-estimates}
\renewcommand{\arraystretch}{1.05}
\begin{tabular}{llllllll}
\toprule
Design & Snapshot & \makecell[b]{Treated\\ available} & Matched & \makecell[b]{Unique\\ controls} & Estimate & 95\% CI & \makecell[b]{Max\\ $|$SMD$|$\\ after} \\
\midrule
Historical & 2016 & 3,825 & 3,825 & 2,749 & 10.17 & [9.66, 10.67] & 0.031 \\
Historical & 2017 & 4,133 & 4,133 & 2,961 & 11.18 & [10.67, 11.69] & 0.026 \\
Prior-attention & 2017 & 3,798 & 3,798 & 2,307 & 1.95 & [1.70, 2.21] & 0.046 \\
Historical & 2022 & 4,373 & 4,373 & 3,197 & 11.26 & [10.71, 11.80] & 0.024 \\
Prior-attention & 2022 & 3,950 & 3,950 & 2,469 & 2.23 & [1.90, 2.56] & 0.041 \\
Historical & 2025 & 4,616 & 4,616 & 3,426 & 11.68 & [11.15, 12.20] & 0.023 \\
Prior-attention & 2025 & 4,349 & 4,349 & 2,829 & 1.96 & [1.69, 2.22] & 0.039 \\
\bottomrule
\end{tabular}
\begin{flushleft}\footnotesize
Nearest-neighbor matching with replacement. Historical designs balance chart year, age, Billboard weeks, peak position, and artist-level historical covariates; prior-attention designs additionally balance the previous high-coverage snapshot's popularity with a five-point caliper. Intervals use paired standard errors; replacement-aware inference is reported in the matching-inference table.
\end{flushleft}
\end{table}


{\footnotesize
\setlength{\tabcolsep}{2.5pt}
\renewcommand{\arraystretch}{1.05}
\begin{longtable}{>{\raggedright\arraybackslash}p{0.08\linewidth}>{\raggedright\arraybackslash}p{0.12\linewidth}>{\raggedright\arraybackslash}p{0.08\linewidth}>{\raggedright\arraybackslash}p{0.17\linewidth}>{\centering\arraybackslash}p{0.068\linewidth}>{\centering\arraybackslash}p{0.068\linewidth}>{\centering\arraybackslash}p{0.068\linewidth}>{\centering\arraybackslash}p{0.068\linewidth}>{\centering\arraybackslash}p{0.068\linewidth}>{\centering\arraybackslash}p{0.068\linewidth}}
\caption{Covariate balance across song-level, strict-temporal, and prior-attention matched designs.}\label{tab:si-matching-balance}\\
\toprule
 &  &  &  & \multicolumn{3}{c}{Before matching} & \multicolumn{3}{c}{After matching} \\
\cmidrule(lr){5-7}\cmidrule(lr){8-10}
Family & Design & Snapshot & Covariate & T mean & C mean & SMD & T mean & C mean & SMD \\
\midrule
\endfirsthead
\multicolumn{10}{@{}l}{\footnotesize Table~\thetable{} continued.}\\
\toprule
 &  &  &  & \multicolumn{3}{c}{Before matching} & \multicolumn{3}{c}{After matching} \\
\cmidrule(lr){5-7}\cmidrule(lr){8-10}
Family & Design & Snapshot & Covariate & T mean & C mean & SMD & T mean & C mean & SMD \\
\midrule
\endhead
\bottomrule
\endfoot
song & Hist. & 2016 & Chart debut year & 1981.64 & 1986.63 & -0.295 & 1981.64 & 1981.84 & -0.013 \\
song & Hist. & 2016 & Song age & 35.16 & 30.17 & 0.295 & 35.16 & 34.96 & 0.013 \\
song & Hist. & 2016 & Billboard weeks & 15.80 & 11.62 & 0.497 & 15.80 & 15.52 & 0.031 \\
song & Hist. & 2016 & Peak chart position & 26.44 & 46.16 & -0.692 & 26.44 & 26.65 & -0.008 \\
song & Hist. & 2016 & Artist catalog size, excluding focal song & 14.70 & 14.32 & 0.019 & 14.70 & 14.59 & 0.007 \\
song & Hist. & 2016 & Artist total Billboard weeks, excluding focal song & 167.35 & 142.43 & 0.143 & 167.35 & 166.20 & 0.007 \\
song & Hist. & 2016 & Artist mean Billboard weeks, excluding focal song & 11.55 & 11.62 & -0.015 & 11.55 & 11.56 & -0.002 \\
song & Hist. & 2016 & Artist best other-song peak position & 13.50 & 14.09 & -0.028 & 13.50 & 13.25 & 0.011 \\
song & Hist. & 2016 & Artist superstar indicator, top 1 percent & 0.17 & 0.14 & 0.075 & 0.17 & 0.17 & 0.000 \\
song & Hist. & 2017 & Chart debut year & 1981.62 & 1984.21 & -0.155 & 1981.62 & 1981.75 & -0.009 \\
song & Hist. & 2017 & Song age & 35.93 & 33.34 & 0.155 & 35.93 & 35.80 & 0.009 \\
song & Hist. & 2017 & Billboard weeks & 15.61 & 10.80 & 0.587 & 15.61 & 15.38 & 0.026 \\
song & Hist. & 2017 & Peak chart position & 27.12 & 48.26 & -0.742 & 27.12 & 27.36 & -0.009 \\
song & Hist. & 2017 & Artist catalog size, excluding focal song & 15.14 & 16.54 & -0.060 & 15.14 & 15.02 & 0.006 \\
song & Hist. & 2017 & Artist total Billboard weeks, excluding focal song & 169.77 & 161.17 & 0.045 & 169.77 & 168.89 & 0.005 \\
song & Hist. & 2017 & Artist mean Billboard weeks, excluding focal song & 11.53 & 11.48 & 0.010 & 11.53 & 11.54 & -0.003 \\
song & Hist. & 2017 & Artist best other-song peak position & 13.37 & 13.82 & -0.021 & 13.37 & 13.16 & 0.010 \\
song & Hist. & 2017 & Artist superstar indicator, top 1 percent & 0.17 & 0.15 & 0.040 & 0.17 & 0.17 & 0.000 \\
song & Prior & 2017 & Chart debut year & 1981.64 & 1985.18 & -0.213 & 1981.64 & 1982.35 & -0.046 \\
song & Prior & 2017 & Song age & 35.91 & 32.37 & 0.213 & 35.91 & 35.20 & 0.046 \\
song & Prior & 2017 & Billboard weeks & 15.80 & 11.30 & 0.547 & 15.80 & 15.55 & 0.029 \\
song & Prior & 2017 & Peak chart position & 26.39 & 46.79 & -0.720 & 26.39 & 26.02 & 0.014 \\
song & Prior & 2017 & Artist catalog size, excluding focal song & 14.71 & 16.21 & -0.073 & 14.71 & 14.58 & 0.008 \\
song & Prior & 2017 & Artist total Billboard weeks, excluding focal song & 167.67 & 162.90 & 0.027 & 167.67 & 166.97 & 0.004 \\
song & Prior & 2017 & Artist mean Billboard weeks, excluding focal song & 11.55 & 11.72 & -0.034 & 11.55 & 11.60 & -0.010 \\
song & Prior & 2017 & Artist best other-song peak position & 13.46 & 13.55 & -0.004 & 13.46 & 12.90 & 0.026 \\
song & Prior & 2017 & Artist superstar indicator, top 1 percent & 0.17 & 0.16 & 0.020 & 0.17 & 0.17 & 0.001 \\
song & Prior & 2017 & Prior Spotify popularity & 42.84 & 29.45 & 0.772 & 42.84 & 42.26 & 0.036 \\
song & Hist. & 2022 & Chart debut year & 1982.57 & 1987.08 & -0.249 & 1982.57 & 1982.78 & -0.013 \\
song & Hist. & 2022 & Song age & 40.08 & 35.57 & 0.249 & 40.08 & 39.87 & 0.013 \\
song & Hist. & 2022 & Billboard weeks & 15.68 & 10.52 & 0.611 & 15.68 & 15.47 & 0.024 \\
song & Hist. & 2022 & Peak chart position & 27.47 & 48.88 & -0.749 & 27.47 & 27.69 & -0.008 \\
song & Hist. & 2022 & Artist catalog size, excluding focal song & 15.01 & 17.84 & -0.110 & 15.01 & 14.89 & 0.006 \\
song & Hist. & 2022 & Artist total Billboard weeks, excluding focal song & 167.52 & 163.81 & 0.018 & 167.52 & 166.49 & 0.005 \\
song & Hist. & 2022 & Artist mean Billboard weeks, excluding focal song & 11.62 & 11.19 & 0.083 & 11.62 & 11.62 & 0.000 \\
song & Hist. & 2022 & Artist best other-song peak position & 13.43 & 14.06 & -0.030 & 13.43 & 13.23 & 0.010 \\
song & Hist. & 2022 & Artist superstar indicator, top 1 percent & 0.16 & 0.16 & 0.022 & 0.16 & 0.16 & 0.000 \\
song & Prior & 2022 & Chart debut year & 1981.79 & 1984.53 & -0.163 & 1981.79 & 1982.42 & -0.041 \\
song & Prior & 2022 & Song age & 40.86 & 38.12 & 0.163 & 40.86 & 40.23 & 0.041 \\
song & Prior & 2022 & Billboard weeks & 15.67 & 10.88 & 0.582 & 15.67 & 15.40 & 0.032 \\
song & Prior & 2022 & Peak chart position & 27.07 & 48.07 & -0.738 & 27.07 & 26.91 & 0.006 \\
song & Prior & 2022 & Artist catalog size, excluding focal song & 15.13 & 16.74 & -0.068 & 15.13 & 14.87 & 0.014 \\
song & Prior & 2022 & Artist total Billboard weeks, excluding focal song & 170.18 & 163.26 & 0.036 & 170.18 & 168.67 & 0.008 \\
song & Prior & 2022 & Artist mean Billboard weeks, excluding focal song & 11.56 & 11.54 & 0.003 & 11.56 & 11.62 & -0.013 \\
song & Prior & 2022 & Artist best other-song peak position & 13.32 & 13.75 & -0.021 & 13.32 & 12.75 & 0.027 \\
song & Prior & 2022 & Artist superstar indicator, top 1 percent & 0.17 & 0.16 & 0.033 & 0.17 & 0.17 & 0.000 \\
song & Prior & 2022 & Prior Spotify popularity & 47.33 & 32.09 & 0.837 & 47.33 & 46.71 & 0.038 \\
song & Hist. & 2025 & Chart debut year & 1982.53 & 1989.45 & -0.363 & 1982.53 & 1982.76 & -0.015 \\
song & Hist. & 2025 & Song age & 43.12 & 36.20 & 0.363 & 43.12 & 42.89 & 0.015 \\
song & Hist. & 2025 & Billboard weeks & 15.57 & 10.12 & 0.638 & 15.57 & 15.37 & 0.023 \\
song & Hist. & 2025 & Peak chart position & 27.69 & 49.44 & -0.760 & 27.69 & 27.90 & -0.007 \\
song & Hist. & 2025 & Artist catalog size, excluding focal song & 14.99 & 18.76 & -0.139 & 14.99 & 14.87 & 0.006 \\
song & Hist. & 2025 & Artist total Billboard weeks, excluding focal song & 166.93 & 167.39 & -0.002 & 166.93 & 165.92 & 0.005 \\
song & Hist. & 2025 & Artist mean Billboard weeks, excluding focal song & 11.61 & 10.99 & 0.119 & 11.61 & 11.61 & 0.001 \\
song & Hist. & 2025 & Artist best other-song peak position & 13.51 & 14.15 & -0.030 & 13.51 & 13.34 & 0.008 \\
song & Hist. & 2025 & Artist superstar indicator, top 1 percent & 0.16 & 0.16 & 0.012 & 0.16 & 0.16 & 0.000 \\
song & Prior & 2025 & Chart debut year & 1982.61 & 1987.07 & -0.247 & 1982.61 & 1983.24 & -0.039 \\
song & Prior & 2025 & Song age & 43.04 & 38.58 & 0.247 & 43.04 & 42.41 & 0.039 \\
song & Prior & 2025 & Billboard weeks & 15.69 & 10.56 & 0.609 & 15.69 & 15.48 & 0.024 \\
song & Prior & 2025 & Peak chart position & 27.42 & 48.70 & -0.745 & 27.42 & 27.16 & 0.009 \\
song & Prior & 2025 & Artist catalog size, excluding focal song & 15.05 & 18.20 & -0.122 & 15.05 & 14.83 & 0.011 \\
song & Prior & 2025 & Artist total Billboard weeks, excluding focal song & 168.00 & 167.25 & 0.004 & 168.00 & 166.56 & 0.008 \\
song & Prior & 2025 & Artist mean Billboard weeks, excluding focal song & 11.62 & 11.24 & 0.075 & 11.62 & 11.65 & -0.005 \\
song & Prior & 2025 & Artist best other-song peak position & 13.46 & 13.77 & -0.015 & 13.46 & 13.15 & 0.015 \\
song & Prior & 2025 & Artist superstar indicator, top 1 percent & 0.17 & 0.16 & 0.015 & 0.17 & 0.17 & 0.000 \\
song & Prior & 2025 & Prior Spotify popularity & 53.20 & 37.43 & 0.826 & 53.20 & 52.59 & 0.037 \\
strict temporal & Exact baseline & 2017 & Baseline Spotify attention (exact) & -- & -- & -- & -- & -- & 0.000 \\
strict temporal & Exact baseline & 2017 & Chart debut year & -- & -- & -- & -- & -- & -0.055 \\
strict temporal & Exact baseline & 2017 & Billboard weeks & -- & -- & -- & -- & -- & 0.012 \\
strict temporal & Exact baseline & 2017 & Peak chart position & -- & -- & -- & -- & -- & -0.036 \\
strict temporal & Exact baseline & 2022 & Baseline Spotify attention (exact) & -- & -- & -- & -- & -- & 0.000 \\
strict temporal & Exact baseline & 2022 & Chart debut year & -- & -- & -- & -- & -- & -0.099 \\
strict temporal & Exact baseline & 2022 & Billboard weeks & -- & -- & -- & -- & -- & 0.048 \\
strict temporal & Exact baseline & 2022 & Peak chart position & -- & -- & -- & -- & -- & -0.034 \\
Baseline-attention placebo & Low baseline (Q1) & pooled & Baseline Spotify popularity & -- & -- & -- & 23.43 & 23.28 & 0.017 \\
Baseline-attention placebo & Low baseline (Q1) & pooled & Baseline song age & -- & -- & -- & 48.87 & 48.72 & 0.019 \\
Baseline-attention placebo & Low baseline (Q1) & pooled & Chart debut year & -- & -- & -- & 1968.57 & 1968.68 & -0.015 \\
Baseline-attention placebo & Low baseline (Q1) & pooled & Billboard weeks & -- & -- & -- & 9.93 & 9.82 & 0.022 \\
Baseline-attention placebo & Low baseline (Q1) & pooled & Peak chart position & -- & -- & -- & 39.51 & 38.76 & 0.026 \\
Baseline-attention placebo & Low baseline (Q1) & pooled & Artist catalog size, excluding focal song & -- & -- & -- & 13.07 & 12.89 & 0.013 \\
Baseline-attention placebo & Low baseline (Q1) & pooled & Artist total Billboard weeks, excluding focal song & -- & -- & -- & 123.40 & 118.95 & 0.029 \\
Baseline-attention placebo & Low baseline (Q1) & pooled & Artist superstar indicator, top 1 percent & -- & -- & -- & 0.11 & 0.11 & 0.000 \\
Baseline-attention placebo & Low baseline (T1) & pooled & Baseline Spotify popularity & -- & -- & -- & 26.74 & 26.70 & 0.004 \\
Baseline-attention placebo & Low baseline (T1) & pooled & Baseline song age & -- & -- & -- & 46.68 & 46.35 & 0.034 \\
Baseline-attention placebo & Low baseline (T1) & pooled & Chart debut year & -- & -- & -- & 1970.76 & 1971.05 & -0.030 \\
Baseline-attention placebo & Low baseline (T1) & pooled & Billboard weeks & -- & -- & -- & 10.70 & 10.65 & 0.010 \\
Baseline-attention placebo & Low baseline (T1) & pooled & Peak chart position & -- & -- & -- & 36.02 & 35.18 & 0.029 \\
Baseline-attention placebo & Low baseline (T1) & pooled & Artist catalog size, excluding focal song & -- & -- & -- & 15.10 & 14.48 & 0.039 \\
Baseline-attention placebo & Low baseline (T1) & pooled & Artist total Billboard weeks, excluding focal song & -- & -- & -- & 151.89 & 148.69 & 0.018 \\
Baseline-attention placebo & Low baseline (T1) & pooled & Artist superstar indicator, top 1 percent & -- & -- & -- & 0.15 & 0.15 & 0.000 \\
Baseline-attention placebo & Regression residual (T1) & pooled & Baseline Spotify popularity & -- & -- & -- & 34.76 & 34.49 & 0.015 \\
Baseline-attention placebo & Regression residual (T1) & pooled & Baseline song age & -- & -- & -- & 31.83 & 31.93 & -0.005 \\
Baseline-attention placebo & Regression residual (T1) & pooled & Chart debut year & -- & -- & -- & 1985.55 & 1985.50 & 0.003 \\
Baseline-attention placebo & Regression residual (T1) & pooled & Billboard weeks & -- & -- & -- & 15.74 & 15.56 & 0.018 \\
Baseline-attention placebo & Regression residual (T1) & pooled & Peak chart position & -- & -- & -- & 29.66 & 29.19 & 0.018 \\
Baseline-attention placebo & Regression residual (T1) & pooled & Artist catalog size, excluding focal song & -- & -- & -- & 12.37 & 11.63 & 0.062 \\
Baseline-attention placebo & Regression residual (T1) & pooled & Artist total Billboard weeks, excluding focal song & -- & -- & -- & 154.22 & 149.98 & 0.024 \\
Baseline-attention placebo & Regression residual (T1) & pooled & Artist superstar indicator, top 1 percent & -- & -- & -- & 0.14 & 0.14 & 0.000 \\
Baseline-attention placebo & Hist. successful, low & pooled & Baseline Spotify popularity & -- & -- & -- & 27.79 & 27.74 & 0.009 \\
Baseline-attention placebo & Hist. successful, low & pooled & Baseline song age & -- & -- & -- & 47.46 & 47.44 & 0.001 \\
Baseline-attention placebo & Hist. successful, low & pooled & Chart debut year & -- & -- & -- & 1969.89 & 1969.95 & -0.005 \\
Baseline-attention placebo & Hist. successful, low & pooled & Billboard weeks & -- & -- & -- & 15.95 & 16.05 & -0.035 \\
Baseline-attention placebo & Hist. successful, low & pooled & Peak chart position & -- & -- & -- & 7.42 & 6.68 & 0.118 \\
Baseline-attention placebo & Hist. successful, low & pooled & Artist catalog size, excluding focal song & -- & -- & -- & 15.16 & 13.37 & 0.125 \\
Baseline-attention placebo & Hist. successful, low & pooled & Artist total Billboard weeks, excluding focal song & -- & -- & -- & 148.05 & 145.58 & 0.015 \\
Baseline-attention placebo & Hist. successful, low & pooled & Artist superstar indicator, top 1 percent & -- & -- & -- & 0.11 & 0.11 & 0.000 \\
Baseline-attention placebo & Older, low attention & pooled & Baseline Spotify popularity & -- & -- & -- & 23.15 & 23.01 & 0.015 \\
Baseline-attention placebo & Older, low attention & pooled & Baseline song age & -- & -- & -- & 49.51 & 49.47 & 0.007 \\
Baseline-attention placebo & Older, low attention & pooled & Chart debut year & -- & -- & -- & 1967.93 & 1967.95 & -0.003 \\
Baseline-attention placebo & Older, low attention & pooled & Billboard weeks & -- & -- & -- & 9.66 & 9.62 & 0.010 \\
Baseline-attention placebo & Older, low attention & pooled & Peak chart position & -- & -- & -- & 40.33 & 39.79 & 0.019 \\
Baseline-attention placebo & Older, low attention & pooled & Artist catalog size, excluding focal song & -- & -- & -- & 13.19 & 12.94 & 0.017 \\
Baseline-attention placebo & Older, low attention & pooled & Artist total Billboard weeks, excluding focal song & -- & -- & -- & 123.59 & 118.40 & 0.033 \\
Baseline-attention placebo & Older, low attention & pooled & Artist superstar indicator, top 1 percent & -- & -- & -- & 0.12 & 0.12 & 0.000 \\
\end{longtable}
\begin{flushleft}\footnotesize Hist.\ = historical matching; Prior = prior-attention matching. Before columns compare all treated with all available controls before matching; after columns compare matched samples. Strict temporal rows report post-matching SMDs for the exact-baseline design by cohort (baseline attention is exactly balanced by construction); pre-matching means are not defined for the cohort-eligibility pools. Baseline-attention placebo rows report matched means and SMDs for the single-assignment reduced-form designs. \end{flushleft}
}


Standard paired errors do not account for control reuse under matching with replacement. Table~\ref{tab:si-song-matching-inference} reports, for the August 2025 designs, the naive paired standard error, a bootstrap over matched treated pairs, and a cluster bootstrap over reused control units, together with no-replacement, stricter-caliper, and coarsened-exact-matching design sensitivities. The song-level matched difference remains positive across these inference variants and designs.

\begin{table}[!htbp]
\centering
\footnotesize
\caption{Replacement-aware inference and design sensitivity for the song-level matched premium (August 2025).}
\label{tab:si-song-matching-inference}
\renewcommand{\arraystretch}{1.05}
\setlength{\tabcolsep}{2.5pt}
\begin{tabular}{>{\raggedright\arraybackslash}p{0.135\linewidth}>{\raggedright\arraybackslash}p{0.15\linewidth}>{\centering\arraybackslash}p{0.058\linewidth}>{\centering\arraybackslash}p{0.078\linewidth}>{\centering\arraybackslash}p{0.082\linewidth}>{\centering\arraybackslash}p{0.052\linewidth}>{\centering\arraybackslash}p{0.09\linewidth}>{\centering\arraybackslash}p{0.082\linewidth}>{\centering\arraybackslash}p{0.1\linewidth}}
\toprule
Comparison & Variant & Pairs & \makecell[b]{Control\\ IDs} & Estimate & \makecell[b]{Naive\\ SE} & \makecell[b]{Bootstrap\\ SE\\ (treated)} & \makecell[b]{Cluster\\ SE\\ (control)} & 95\% CI \\
\midrule
baseline-attention + historical + artist-success NN | 2025 (with replacement) & with replacement & 4,349 & 2,829 & 1.96 & 0.135 & 0.128 & 0.167 & [1.63, 2.29] \\
historical + artist-success NN | 2025 (with replacement) & with replacement & 4,616 & 3,426 & 11.68 & 0.268 & 0.263 & 0.333 & [11.02, 12.36] \\
historical + artist-success NN | 2025 (no replacement) & no replacement & 4,561 & 4,561 & 12.53 & 0.278 & -- & -- & [11.99, 13.08] \\
baseline-attention + historical + artist-success NN | 2025 (no replacement) & no replacement & 3,746 & 3,745 & 2.32 & 0.154 & -- & -- & [2.02, 2.62] \\
baseline-attention + historical + artist-success NN | 2025 (stricter caliper 2.0) & stricter caliper & 4,341 & 2,844 & 1.34 & 0.134 & -- & -- & [1.08, 1.61] \\
baseline-attention + historical + artist-success NN | 2025 (CEM decade x baseline-pop quintile) & CEM & 4,339 & -- & 3.03 & -- & -- & -- & -- \\
\bottomrule
\end{tabular}
\begin{flushleft}\footnotesize
Matching with replacement reuses controls. The cluster bootstrap over reused control IDs is the primary uncertainty check, with the treated-pair bootstrap and naive paired SE shown for comparison. No-replacement, stricter-caliper, and coarsened-exact-matching (chart decade $\times$ baseline-popularity quintile) variants are design sensitivities. The interval column reports the cluster-by-control interval where computed and the naive interval otherwise. Control IDs are alphanumeric song-artist identifiers. No-replacement matching excludes repeated record indices, not different records sharing an ID.
\end{flushleft}
\end{table}


{\sloppy The pooled prior-attention matched excess-attention comparison in the main text (2.52 excess-attention points over 12,095 matched pairs) reuses control songs and repeats Billboard song identities across snapshots on both sides of the comparison. Table~\ref{tab:si-matched-excess-inference} therefore reports replacement-aware uncertainty for this estimand on the fixed pair set: the naive paired standard error, a bootstrap over matched pairs, a block bootstrap over treated song records, a cluster bootstrap over reused control song records, and a two-way cluster bootstrap that independently resamples treated and control song records and reweights pairs by the product of their cluster draw counts. The point estimate is identical across rows by construction. Only the uncertainty calculation varies. The two-way interval is primary because it accounts jointly for both dependence structures, and the difference remains positive under it.\par}

\begin{table}[!htbp]
\centering
\footnotesize
\caption{Replacement-aware inference for the pooled prior-attention matched-excess premium.}
\label{tab:si-matched-excess-inference}
\renewcommand{\arraystretch}{1.05}
\setlength{\tabcolsep}{3pt}
\begin{tabular}{>{\raggedright\arraybackslash}p{0.26\linewidth}>{\centering\arraybackslash}p{0.082\linewidth}>{\centering\arraybackslash}p{0.05\linewidth}>{\centering\arraybackslash}p{0.09\linewidth}>{\centering\arraybackslash}p{0.065\linewidth}>{\centering\arraybackslash}p{0.072\linewidth}>{\centering\arraybackslash}p{0.072\linewidth}>{\raggedright\arraybackslash}p{0.09\linewidth}}
\toprule
Method & Estimate & SE & 95\% CI & Pairs & \makecell[b]{Treated\\ songs} & \makecell[b]{Control\\ songs} & Role \\
\midrule
naive paired SE & 2.52 & 0.088 & [2.35, 2.69] & 12,095 & 4,534 & 4,764 &  \\
bootstrap over matched pairs & 2.52 & 0.087 & [2.35, 2.69] & 12,095 & 4,534 & 4,764 &  \\
block bootstrap over treated song records & 2.52 & 0.089 & [2.34, 2.69] & 12,095 & 4,534 & 4,764 &  \\
cluster bootstrap over reused control song records & 2.52 & 0.128 & [2.27, 2.78] & 12,095 & 4,534 & 4,764 &  \\
two-way pigeonhole bootstrap (treated x control song records) & 2.52 & 0.181 & [2.16, 2.87] & 12,095 & 4,534 & 4,764 & primary \\
\bottomrule
\end{tabular}
\begin{flushleft}\footnotesize
All methods evaluate the identical estimand on the identical 12{,}095-pair set with fixed pair-level excess differences; only the uncertainty calculation varies. The two-way pigeonhole bootstrap independently resamples treated song records and reused control song records and reweights pairs by the product of their cluster draw counts; it is primary because the pooled comparison both reuses control songs and repeats Billboard song identities across snapshots. Percentile intervals use 2{,}000 replications.
\end{flushleft}
\end{table}


\section*{Pre-Film Attention and Integrated Inference}

The pre-film attention subanalysis selects film-linked songs in the bottom tercile of pre-film excess attention among the 355 film-linked songs with an observed Spotify snapshot strictly before and at or after their first film year. The 119 selected songs are identified by exact Billboard record identity. Baseline is the latest observed snapshot strictly before the first film year. The resulting threshold is +9.876 excess-attention points, so membership does not require attention below the expected curve or low absolute Spotify popularity. Of the 119 songs, 77 have nonnegative pre-film excess attention. The lower-tercile group is therefore defined relative to excess attention within the eligible film-linked group, without a requirement of prior decline or a minimum age. The above-threshold outcome indicates whether a song exceeds this fixed threshold at the first observed post-film snapshot. The sparse snapshots and year-level film dates do not establish attention immediately before or after release. Never-film controls receive pseudo-event years drawn from the empirical distribution of treated first-film years, pass through the identical pre/post rule and treated-derived baseline threshold, and are matched with replacement on baseline excess attention, baseline popularity, baseline age, chart year, Billboard weeks, peak position, and artist-level historical covariates. The contrast is the pre-to-post change in excess attention among selected film-linked songs minus the corresponding change among their matched controls.

Table~\ref{tab:si-dormant-integrated-inference} reports the integrated inference from 500 deterministic-seed pseudo-event assignments and 250 full nested-bootstrap replications that resample never-film song trajectories, refit the snapshot-specific MemoryDecay curves, recompute expected and excess attention, reconstruct the lower tercile within the eligible film-linked sample, resample the selected treated records, reassign pseudo-events to controls, and rematch controls inside every replication. Bootstrap curves use the same estimator as the point-estimate baselines: floor-age cells with mean fractional age within each cell, positive age-cell attention means, a log-response squared-error objective, capped inverse-age weights, the same log/logit parameterization, and the full multi-start grid. All 500 assignment-specific contrasts are positive (Figure~\ref{fig:si-dormant-pseudoevent}). The assignment interval describes design stability rather than sampling uncertainty, which is carried by the integrated bootstrap. The assignments are repeated realizations of the same comparison, not independent empirical replications.

\begin{table}[!htbp]
\centering
\footnotesize
\caption{Lower pre-film excess attention: support and Monte Carlo-integrated inference.}
\label{tab:si-dormant-integrated-inference}
\renewcommand{\arraystretch}{1.05}
\begin{tabular}{>{\raggedright\arraybackslash}p{0.52\linewidth}>{\centering\arraybackslash}p{0.18\linewidth}}
\toprule
Quantity & Value \\
\midrule
film-linked songs with pre-film and post-film Spotify snapshots & 355 \\
bottom-tercile pre-film excess-attention threshold (q33) & 9.876 \\
lower-tercile pre-film excess-attention treated songs & 119 \\
Pseudo-event assignments & 500 \\
Assignments with positive contrast & 500 \\
Assignment mean (integrated point estimate) & 5.044 \\
Assignment median & 5.029 \\
Assignment SD & 0.435 \\
Assignment interval, 2.5th pct & 4.243 \\
Assignment interval, 97.5th pct & 5.997 \\
Integrated bootstrap 95\% CI, low & 2.972 \\
Integrated bootstrap 95\% CI, high & 6.868 \\
Valid nested-bootstrap replications & 250 \\
Treated above-threshold probability & 0.395 \\
Matched-control above-threshold probability & 0.215 \\
Above-threshold probability difference (pp) & 17.987 \\
Above-threshold integrated 95\% CI, low (pp) & 6.736 \\
Above-threshold integrated 95\% CI, high (pp) & 27.593 \\
Decay-equivalent years (median) & 6.883 \\
Decay-equivalent 95\% CI, low & 3.890 \\
Decay-equivalent 95\% CI, high & 10.225 \\
Inversion failure rate & 0.000 \\
Median 95th-pct control reuse & 3 \\
Maximum control reuse & 7 \\
Matched lower-tercile treated songs per assignment & 119 \\
\bottomrule
\end{tabular}
\begin{flushleft}\footnotesize
Inference integrates 500 pseudo-event assignments and 250 nested-bootstrap replications with snapshot-specific MemoryDecay refitting, excess-attention reconstruction, pseudo-event reassignment, and rematching. The assignment interval is a design-stability distribution across pseudo-event assignments, not a confidence interval. The analytical support comprises 355 film-linked songs with pre-film and post-film snapshots; the 119 songs in the lower tercile of pre-film excess attention are matched in every pseudo-event assignment. Bootstrap curves use the same floor-age aggregation, mean age within each cell, log-response squared-error objective, capped inverse-age weights, parameter constraints, and full starting grid as the point-estimate curves.
\end{flushleft}
\end{table}


\begin{figure}[H]
\centering
\includegraphics[width=.8\linewidth]{Figures/SI/figure_si_dormant_pseudoevent.pdf}
\caption{Pseudo-event assignment stability of the lower pre-film excess-attention contrast. The histogram shows the 500 assignment-specific placebo-adjusted excess-attention change contrasts. The solid line marks the integrated mean and the dotted lines the 2.5th--97.5th percentile assignment band. Every assignment yields a positive contrast. The band describes assignment stability, not a confidence interval.}
\label{fig:si-dormant-pseudoevent}
\end{figure}

Low baseline attention can mechanically predict later increases, so change is also evaluated under alternative raw-popularity definitions with a single seeded pseudo-event assignment (Table~\ref{tab:si-dormant-definitions}): bottom-quartile and bottom-tercile baseline popularity, the bottom tercile of residuals from a regression of baseline popularity on age, squared age, weeks, and peak position, historically successful low-attention songs, and older low-attention songs. Across the four larger-sample definitions the placebo-adjusted change is positive with bootstrap intervals excluding zero and permutation $p$-values below 0.001. The historically successful low-attention definition has low support and is exploratory. These checks assess sensitivity to baseline classification. The regression-residual classification differs from the primary MemoryDecay excess-attention classification. These comparisons do not test whether adjusted changes differ across the primary excess-attention terciles.

\begin{table}[!htbp]
\centering
\footnotesize
\caption{Baseline-attention definitions: support and single-assignment placebo-adjusted changes.}
\label{tab:si-dormant-definitions}
\renewcommand{\arraystretch}{1.05}
\setlength{\tabcolsep}{2.5pt}
\begin{tabular}{>{\raggedright\arraybackslash}p{0.13\linewidth}>{\centering\arraybackslash}p{0.088\linewidth}>{\centering\arraybackslash}p{0.055\linewidth}>{\centering\arraybackslash}p{0.072\linewidth}>{\centering\arraybackslash}p{0.075\linewidth}>{\centering\arraybackslash}p{0.078\linewidth}>{\centering\arraybackslash}p{0.108\linewidth}>{\centering\arraybackslash}p{0.068\linewidth}>{\centering\arraybackslash}p{0.108\linewidth}>{\centering\arraybackslash}p{0.065\linewidth}}
\toprule
Definition & \makecell[b]{Selected\\ treated} & Pairs & \makecell[b]{Treated\\ change} & \makecell[b]{Matched\\ control\\ change} & \makecell[b]{Placebo-\\ adjusted} & \makecell[b]{Bootstrap\\ 95\% CI} & \makecell[b]{Perm.\\ $p$} & \makecell[b]{No-\\replacement} & \makecell[b]{Max\\ $|$SMD$|$} \\
\midrule
Low baseline Spotify attention, bottom quartile & 97 & 97 & 10.93 & 3.87 & 7.06 & [4.65, 9.43] & $<0.001$ & 6.71 & 0.029 \\
Low baseline Spotify attention, bottom tercile & 124 & 124 & 10.16 & 4.60 & 5.57 & [3.56, 7.76] & $<0.001$ & 5.29 & 0.039 \\
Baseline-popularity regression residual, bottom tercile & 119 & 119 & 9.39 & 5.64 & 3.76 & [1.92, 5.76] & $<0.001$ & 3.60 & 0.062 \\
Historically successful but low baseline attention & 19 & 19 & 10.84 & 6.11 & 4.74 & [-0.05, 9.74] & 0.737 & 4.74 & 0.125 \\
Older low-attention songs & 94 & 94 & 10.93 & 3.20 & 7.72 & [5.33, 10.22] & $<0.001$ & 7.53 & 0.033 \\
\bottomrule
\end{tabular}
\begin{flushleft}\footnotesize
Four definitions use raw baseline Spotify popularity; the fifth uses the bottom tercile of residuals from a regression of baseline popularity on age, age squared, Billboard weeks, and peak position. Treated-derived thresholds are applied to both treated songs and pseudo-event never-film controls; one seeded pseudo-event assignment (single-assignment robustness, not the primary integrated design). The historically-successful low-attention definition has low support and is exploratory. Bootstrap intervals resample matched treated songs (1{,}000 draws); permutation $p$-values shuffle film-linked labels within the pooled set selected by each definition (2{,}000 draws).
\end{flushleft}
\end{table}


Figure~\ref{fig:si-dormant-rtm} compares the raw attention changes among selected film-linked songs and their matched pseudo-event controls under each definition. Both groups show increases, while a positive difference remains in the four larger-sample comparisons.

\begin{figure}[H]
\centering
\includegraphics[width=.85\linewidth]{Figures/SI/figure_si_dormant_rtm.pdf}
\caption{Attention changes across alternative pre-film attention definitions. Gray squares are raw first-post-event changes for selected film-linked songs. Open circles are the corresponding matched pseudo-event control changes. Filled points are the differences in change with bootstrap 95\% intervals from single-assignment comparisons. Both groups show attention increases, with a positive difference across the four larger-sample definitions. The historically successful low-attention definition has low support and is exploratory (lighter interval).}
\label{fig:si-dormant-rtm}
\end{figure}

\section*{Repeated Embedding and Film Visibility}

Repeated embedding is measured with movie-count categories. The primary Spotify specification counts only films released by each snapshot year and clusters inference by song across repeated song-snapshot observations. Table~\ref{tab:si-repeated-embedding} reports this gradient, the comparison using lifetime film counts, and the Last.fm cumulative listener-reach gradient (lifetime categories, single snapshot). The positive Spotify gradient persists when film counts exclude films with release years later than the snapshot year. These gradients describe adjusted associations and do not distinguish selection from changes following placement.

\begin{table}[!htbp]
\centering
\footnotesize
\caption{Repeated film embedding: time-respecting primary specification, lifetime comparison, and Last.fm.}
\label{tab:si-repeated-embedding}
\renewcommand{\arraystretch}{1.05}
\setlength{\tabcolsep}{3pt}
\begin{tabular}{>{\raggedright\arraybackslash}p{0.25\linewidth}>{\raggedright\arraybackslash}p{0.17\linewidth}>{\centering\arraybackslash}p{0.082\linewidth}>{\centering\arraybackslash}p{0.05\linewidth}>{\centering\arraybackslash}p{0.10\linewidth}>{\centering\arraybackslash}p{0.06\linewidth}>{\centering\arraybackslash}p{0.06\linewidth}}
\toprule
Model & Term & Estimate & SE & 95\% CI & N obs. & N songs \\
\midrule
Repeated embedding (time-respecting, primary) & 1 realized films vs none & 9.25 & 0.264 & [8.73, 9.76] & 83,780 & 28,281 \\
Repeated embedding (time-respecting, primary) & 2-3 realized films vs none & 16.72 & 0.341 & [16.05, 17.39] & 83,780 & 28,281 \\
Repeated embedding (time-respecting, primary) & 4+ realized films vs none & 24.19 & 0.424 & [23.35, 25.02] & 83,780 & 28,281 \\
Repeated embedding (lifetime comparison) & 1 lifetime films vs none & 9.28 & 0.268 & [8.75, 9.80] & 83,780 & 28,281 \\
Repeated embedding (lifetime comparison) & 2-3 lifetime films vs none & 16.43 & 0.348 & [15.75, 17.11] & 83,780 & 28,281 \\
Repeated embedding (lifetime comparison) & 4+ lifetime films vs none & 24.08 & 0.420 & [23.26, 24.91] & 83,780 & 28,281 \\
Last.fm repeated embedding (lifetime) & 1 lifetime films vs none & 1.34 & 0.038 & [1.26, 1.41] & 17,988 & -- \\
Last.fm repeated embedding (lifetime) & 2-3 lifetime films vs none & 2.06 & 0.047 & [1.97, 2.15] & 17,988 & -- \\
Last.fm repeated embedding (lifetime) & 4+ lifetime films vs none & 2.78 & 0.053 & [2.67, 2.88] & 17,988 & -- \\
\bottomrule
\end{tabular}
\begin{flushleft}\footnotesize
Spotify models adjust for age, squared age, Billboard weeks, peak position, and snapshot fixed effects; the primary specification counts only film appearances realized by each snapshot and reports cluster-by-song intervals. The lifetime comparison shows the gradient is not an artifact of attaching later films to earlier snapshots. Last.fm uses log1p cumulative listeners with lifetime counts (single snapshot) and HC1 intervals.
\end{flushleft}
\end{table}


Figure~\ref{fig:si-crossplatform-embedding} shows monotonic gradients on both platforms. Spotify film counts use release-year eligibility. Last.fm has one cumulative listener snapshot and uses lifetime movie-count categories. Its panel describes a complementary listener-reach association rather than a temporal replication.

\begin{figure}[H]
\centering
\includegraphics[width=.95\linewidth]{Figures/SI/figure_si_crossplatform_embedding.pdf}
\caption{Film-count gradients in Spotify popularity and cumulative Last.fm listener reach. Panel A: pooled Spotify differences by film count through the snapshot year (cluster-by-song 95\% CIs). Panel B: cumulative Last.fm listener-reach differences by lifetime movie-count category in the July 2017 snapshot (HC1 95\% CIs). Both platforms show monotonic associations under their respective film-count definitions.}
\label{fig:si-crossplatform-embedding}
\end{figure}

{\sloppy Film visibility uses recorded worldwide box office for films released by the snapshot year. Table~\ref{tab:si-visibility-models} reports the continuous model behind the main visibility estimate. Table~\ref{tab:si-visibility-selection} splits eligible song-snapshots with at least one film appearance by the availability of box-office data for films released by that year. The row with observed box office is the model's estimation sample. The 4,272 records with box-office metadata in the global film-linked population and the 4,294 song identifiers in the visibility model use different linkage and counting conventions. The visibility model draws on snapshot-specific alphanumeric matches, including 57 Billboard records outside the global film-linked classification; 32 records in the global coverage count lack an eligible Spotify observation. Its 4,297 distinct Billboard records correspond to 4,294 alphanumeric song identifiers. The manual audit samples the global linkage and does not separately validate the additional alphanumeric matches. In these tables, realized exposure refers to film release-year eligibility. The recorded box-office totals may include later receipts. Box office measures film-level visibility, not the prominence of a song within a scene or soundtrack.\par}

\begin{table}[!htbp]
\centering
\footnotesize
\caption{Continuous time-respecting film-visibility model.}
\label{tab:si-visibility-models}
\renewcommand{\arraystretch}{1.05}
\setlength{\tabcolsep}{3pt}
\begin{tabular}{>{\raggedright\arraybackslash}p{0.22\linewidth}>{\raggedright\arraybackslash}p{0.20\linewidth}>{\centering\arraybackslash}p{0.082\linewidth}>{\centering\arraybackslash}p{0.07\linewidth}>{\centering\arraybackslash}p{0.095\linewidth}>{\centering\arraybackslash}p{0.06\linewidth}>{\centering\arraybackslash}p{0.06\linewidth}}
\toprule
Model & Term & Estimate & \makecell[b]{Cluster\\ SE} & 95\% CI & N obs. & N songs \\
\midrule
Film visibility (continuous, time-respecting, primary) & log1p max recorded worldwide box office; films released by snapshot year & 1.15 & 0.102 & [0.95, 1.35] & 15,245 & 4,294 \\
\bottomrule
\end{tabular}
\begin{flushleft}\footnotesize
Song-snapshots linked to films released by the snapshot year with recorded worldwide box office. Recorded totals may include receipts after the snapshot date. The model adjusts for age, squared age, Billboard weeks, peak position, realized film count, and snapshot fixed effects, with cluster-by-song intervals. Box office measures film-level visibility, not scene-level soundtrack prominence.
\end{flushleft}
\end{table}


\begin{table}[!htbp]
\centering
\footnotesize
\caption{Support and selection diagnostics for the film-visibility model.}
\label{tab:si-visibility-selection}
\renewcommand{\arraystretch}{1.05}
\setlength{\tabcolsep}{3pt}
\begin{tabular}{>{\raggedright\arraybackslash}p{0.21\linewidth}>{\centering\arraybackslash}p{0.06\linewidth}>{\centering\arraybackslash}p{0.055\linewidth}>{\centering\arraybackslash}p{0.05\linewidth}>{\centering\arraybackslash}p{0.055\linewidth}>{\centering\arraybackslash}p{0.05\linewidth}>{\centering\arraybackslash}p{0.085\linewidth}>{\centering\arraybackslash}p{0.10\linewidth}}
\toprule
Time-respecting film-linked observations & N obs. & N songs & \makecell[b]{Mean\\ age} & \makecell[b]{Mean\\ weeks} & \makecell[b]{Mean\\ peak} & \makecell[b]{Mean\\ realized\\ films} & \makecell[b]{Mean\\ popularity} \\
\midrule
without recorded worldwide box office & 1,301 & 430 & 44.3 & 12.9 & 33.9 & 1.06 & 40.9 \\
with recorded worldwide box office & 15,245 & 4,294 & 38.7 & 15.9 & 26.5 & 2.55 & 50.3 \\
\bottomrule
\end{tabular}
\begin{flushleft}\footnotesize
The universe is every song-snapshot linked to at least one film released by the snapshot year. Rows split this universe by whether positive worldwide box office is recorded for at least one eligible film. Recorded totals may include receipts after the snapshot date. The row with observed box office is exactly the estimation sample of the continuous visibility model and its tercile visualization (15,245 song-snapshots, 4,294 normalized song-artist identities); the model's complete-case condition drops no observations.
\end{flushleft}
\end{table}


\section*{Temporal Consistency Checks}

Broad event-time models measure attention relative to the first observed film appearance, with never-film songs and the immediate pre-year forming the omitted reference. Table~\ref{tab:si-temporal-event-time} reports the contrasts with per-bin support. The positive pre-event coefficient documents an attention difference before observed placement, consistent with selection. It does not identify the reasons songs were selected or establish a causal effect of placement.

\begin{table}[!htbp]
\centering
\footnotesize
\caption{Broad song event-time model relative to the first observed film appearance.}
\label{tab:si-temporal-event-time}
\renewcommand{\arraystretch}{1.05}
\begin{tabular}{lllllll}
\toprule
Event bin & Estimate & HC1 SE & HC1 95\% CI & Cluster 95\% CI & N film-linked obs. & N songs \\
\midrule
pre (<= -2 years) & 9.50 & 0.671 & [8.18, 10.81] & [7.80, 11.19] & 343 & 213 \\
event year & 11.23 & 0.908 & [9.45, 13.01] & [9.44, 13.02] & 184 & 184 \\
1-5 years post & 10.16 & 0.357 & [9.46, 10.86] & [9.31, 11.02] & 1,046 & 747 \\
>5 years post & 14.18 & 0.136 & [13.91, 14.45] & [13.71, 14.65] & 15,102 & 4,503 \\
\bottomrule
\end{tabular}
\begin{flushleft}\footnotesize
All 83{,}785 song-snapshot observations; never-film songs and the immediate pre-year form the omitted reference. Adjusts for age, squared age, Billboard weeks, peak position, and snapshot fixed effects. The positive $\leq -2$ bin is direct evidence of selection into film placement. Support columns count film-linked observations in each bin.
\end{flushleft}
\end{table}


The strict exact-pre-attention matched comparison matches treated songs (first film in 2017 or 2022) one-to-one without replacement to never-film controls with identical pre-film Spotify attention and similar chart histories, preserving 93 pairs (60 in the 2017 cohort, 33 in the 2022 cohort). Exact pre-film attention matching yields the same matched support across the evaluated nonnegative caliper settings, so redundant caliper rows are omitted. Table~\ref{tab:si-strict-temporal} reports pooled and cohort-level contrasts with paired $t$-tests and sign-flip randomization inference. The same 93 pairs re-expressed on the excess-attention scale appear in the main text. Pre-film attention is exactly balanced by construction, and post-matching chart-history imbalances are small (Table~\ref{tab:si-matching-balance}).

\begin{table}[!htbp]
\centering
\footnotesize
\caption{Strict exact-pre-attention matched temporal comparison (raw Spotify scale).}
\label{tab:si-strict-temporal}
\renewcommand{\arraystretch}{1.05}
\begin{tabular}{lllllll}
\toprule
Level & Contrast & Pairs & Estimate & SE & $t$-test $p$ & Randomization $p$ \\
\midrule
pooled & event year & 93 & 2.43 & 0.936 & 0.011 & 0.010 \\
pooled & first post & 93 & 2.94 & 0.919 & 0.002 & 0.001 \\
2017 cohort & event year & 60 & 2.30 & 0.890 & 0.012 & 0.011 \\
2017 cohort & first post & 60 & 3.12 & 1.153 & 0.009 & 0.010 \\
2017 cohort & second post & 60 & 3.05 & 1.452 & 0.040 & 0.044 \\
2022 cohort & event year & 33 & 2.67 & 2.107 & 0.215 & 0.240 \\
2022 cohort & first post & 33 & 2.61 & 1.542 & 0.101 & 0.105 \\
2022 cohort & pre-period placebo & 33 & 0.27 & 1.019 & 0.791 & 0.818 \\
\bottomrule
\end{tabular}
\begin{flushleft}\footnotesize
Optimal 1:1 matching without replacement under a hard zero-point caliper on baseline Spotify attention with standardized chart-history distance as tie-breaker. The 2017 cohort (60 pairs) uses October 2016 as baseline; the 2022 cohort (33 pairs) uses July 2017; the 2022 pre-placebo contrast checks the 2016--2017 pre-period. Because exact prior-attention matching yields the same matched support across the evaluated nonnegative caliper settings (0, 1, 2, 3, 5 points), redundant caliper rows are omitted. Randomization $p$-values use 5{,}000 sign-flip draws.
\end{flushleft}
\end{table}


\section*{Artist-Catalog Attention Analyses}

Artist catalogs consist of Billboard songs sharing the same normalized artist or band identity. The analyses include only songs with no observed film appearance and do not cover each artist's complete repertoire. Normalized artist/band strings define both the exposure indicator and the artist-level covariates. Static models compare songs by artists with another film-linked Billboard song with songs by artists with no film-linked Billboard song (Table~\ref{tab:si-catalog-static}). Dynamic models measure event time relative to the first film appearance of any other Billboard song by the same artist (Table~\ref{tab:si-catalog-dynamic}). The positive pre-event coefficient shows that an attention difference is present before the first observed film appearance of another song by the artist (Figure~\ref{fig:si-catalog-event-time-selection}). This is consistent with selection and limits interpretation of later differences as changes induced by placement. Only one Last.fm listener snapshot exists, so a Last.fm dynamic model is not estimated.

\begin{table}[!htbp]
\centering
\footnotesize
\caption{Static artist-catalog attention differences among non-film Billboard songs.}
\label{tab:si-catalog-static}
\renewcommand{\arraystretch}{1.05}
\begin{tabular}{llllllll}
\toprule
Platform & Snapshot & Estimate & 95\% CI & $p$ & N obs. & Exposed & Control \\
\midrule
Spotify & 2016 & 2.32 & [1.81, 2.83] & $<0.001$ & 10,966 & 6,858 & 4,108 \\
Spotify & 2017 & 3.57 & [3.05, 4.09] & $<0.001$ & 12,550 & 8,253 & 4,297 \\
Spotify & 2022 & 2.08 & [1.59, 2.57] & $<0.001$ & 15,487 & 9,356 & 6,131 \\
Spotify & 2025 & 2.34 & [1.90, 2.79] & $<0.001$ & 18,442 & 10,671 & 7,771 \\
Spotify & pooled & 2.47 & [2.22, 2.72] & $<0.001$ & 57,445 & 35,138 & 22,307 \\
Last.fm & 2017 & 0.68 & [0.60, 0.75] & $<0.001$ & 12,419 & 8,179 & 4,240 \\
\bottomrule
\end{tabular}
\begin{flushleft}\footnotesize
Exposure: the same normalized artist/band string has another Billboard song linked to a film. Models adjust for song age, squared age, Billboard weeks, peak position, artist catalog size, artist total weeks excluding the focal song, best other-song peak, and superstar status; the pooled model adds snapshot fixed effects (HC1 intervals; the cluster-by-song variant appears in the SI inference table). Last.fm uses log1p cumulative listeners. Estimates describe adjusted attention differences across the observed artist groups.
\end{flushleft}
\end{table}


\begin{table}[!htbp]
\centering
\footnotesize
\caption{Dynamic artist-catalog event-time pattern among non-film songs.}
\label{tab:si-catalog-dynamic}
\renewcommand{\arraystretch}{1.05}
\begin{tabular}{lllllll}
\toprule
Event bin & Estimate & SE & 95\% CI & N exposed obs. & N songs & N artists \\
\midrule
pre (<= -2 years) & 5.31 & 0.512 & [4.31, 6.32] & 407 & 233 & 38 \\
event year & 5.09 & 0.741 & [3.64, 6.54] & 211 & 210 & 43 \\
1-5 years post & 5.40 & 0.254 & [4.90, 5.90] & 1,776 & 1,292 & 171 \\
>5 years post & 3.66 & 0.105 & [3.46, 3.87] & 32,506 & 10,274 & 1,525 \\
\bottomrule
\end{tabular}
\begin{flushleft}\footnotesize
Event time is relative to the first film appearance of any other Billboard song by the same artist; songs by artists with no film-linked song form the omitted reference with the immediate pre-year (66{,}838 observations; HC1 intervals). The positive pre-event estimate documents prior attention differences between groups, so this pattern does not isolate increases associated with placement. Only one Last.fm listener snapshot exists, so a Last.fm dynamic model is not estimated.
\end{flushleft}
\end{table}


\begin{figure}[H]
\centering
\includegraphics[width=.62\linewidth]{Figures/SI/figure_si_catalog_event_time_selection.pdf}
\caption{Artist-catalog attention differences before and after observed film placement. Points are the coefficients of Table~\ref{tab:si-catalog-dynamic}, with HC1 95\% CIs, for the pre, event-year, one-to-five-year post, and more-than-five-year post windows. Event time is relative to the first film appearance of another Billboard song by the same artist. The positive pre-event coefficient documents an attention difference that precedes the observed placement.}
\label{fig:si-catalog-event-time-selection}
\end{figure}

Matched catalog comparisons assess attention differences after balancing song and artist characteristics (Table~\ref{tab:si-catalog-matching}). Table~\ref{tab:si-catalog-balance} reports covariate-level balance because artist catalogs differ in historical visibility and size. Catalog controls are reused heavily under matching with replacement, so Table~\ref{tab:si-catalog-matching-inference} reports cluster-by-control inference and design sensitivities. Under replacement-aware inference, the August 2022 prior-attention contrast is not distinguishable from zero, while no-replacement and coarsened-exact-matching variants return small positive differences. These comparisons estimate conditional associations rather than causal catalog effects.

\begin{table}[H]
\centering
\footnotesize
\caption{Matched artist-catalog attention premiums.}
\label{tab:si-catalog-matching}
\renewcommand{\arraystretch}{1.05}
\setlength{\tabcolsep}{3pt}
\begin{tabular}{>{\raggedright\arraybackslash}p{0.085\linewidth}>{\raggedright\arraybackslash}p{0.09\linewidth}>{\raggedright\arraybackslash}p{0.095\linewidth}>{\centering\arraybackslash}p{0.088\linewidth}>{\centering\arraybackslash}p{0.075\linewidth}>{\centering\arraybackslash}p{0.078\linewidth}>{\centering\arraybackslash}p{0.082\linewidth}>{\centering\arraybackslash}p{0.10\linewidth}>{\centering\arraybackslash}p{0.095\linewidth}}
\toprule
Platform & Design & Snapshot & \makecell[b]{Treated\\ available} & Matched & \makecell[b]{Unique\\ controls} & Estimate & 95\% CI & \makecell[b]{Max\\ $|$SMD$|$\\ after} \\
\midrule
Spotify & Historical & 2016 & 6,858 & 6,858 & 1,997 & 2.49 & [2.13, 2.85] & 0.125 \\
Spotify & Historical & 2017 & 8,253 & 8,253 & 2,201 & 2.89 & [2.54, 3.24] & 0.159 \\
Spotify & Prior-attention & 2017 & 6,794 & 6,794 & 1,927 & 0.46 & [0.29, 0.63] & 0.163 \\
Spotify & Historical & 2022 & 9,356 & 9,356 & 2,932 & 1.38 & [1.02, 1.74] & 0.154 \\
Spotify & Prior-attention & 2022 & 7,712 & 7,712 & 2,197 & -0.31 & [-0.57, -0.05] & 0.201 \\
Spotify & Historical & 2025 & 10,671 & 10,671 & 3,435 & 1.46 & [1.12, 1.79] & 0.178 \\
Spotify & Prior-attention & 2025 & 9,285 & 9,285 & 3,050 & 0.44 & [0.26, 0.62] & 0.168 \\
Spotify & Prior-attention & pooled (inverse-variance) & 23,791 & 23,791 & -- & 0.31 & [0.20, 0.42] & 0.201 \\
Spotify & Historical & pooled (inverse-variance) & 35,138 & 35,138 & -- & 2.04 & [1.87, 2.22] & 0.178 \\
Last.fm & Prior-attention & 2017 & 6,727 & 6,727 & 1,892 & 0.29 & [0.25, 0.33] & 0.165 \\
\bottomrule
\end{tabular}
\begin{flushleft}\footnotesize
Nearest-neighbor matching with replacement among non-film songs. Pooled rows aggregate per-snapshot matched differences by inverse variance (descriptive). Intervals use paired standard errors; replacement-aware inference is reported in the catalog matching-inference table, under which the August 2022 baseline-attention contrast is not robustly distinguishable from zero.
\end{flushleft}
\end{table}


{\footnotesize
\setlength{\tabcolsep}{2.5pt}
\renewcommand{\arraystretch}{1.05}
\begin{longtable}{>{\raggedright\arraybackslash}p{0.078\linewidth}>{\raggedright\arraybackslash}p{0.075\linewidth}>{\raggedright\arraybackslash}p{0.08\linewidth}>{\raggedright\arraybackslash}p{0.185\linewidth}>{\centering\arraybackslash}p{0.068\linewidth}>{\centering\arraybackslash}p{0.068\linewidth}>{\centering\arraybackslash}p{0.068\linewidth}>{\centering\arraybackslash}p{0.068\linewidth}>{\centering\arraybackslash}p{0.068\linewidth}>{\centering\arraybackslash}p{0.068\linewidth}}
\caption{Covariate balance for the artist-catalog matched designs (long format).}\label{tab:si-catalog-balance}\\
\toprule
 &  &  &  & \multicolumn{3}{c}{Before matching} & \multicolumn{3}{c}{After matching} \\
\cmidrule(lr){5-7}\cmidrule(lr){8-10}
Platform & Design & Snapshot & Covariate & T mean & C mean & SMD & T mean & C mean & SMD \\
\midrule
\endfirsthead
\multicolumn{10}{@{}l}{\footnotesize Table~\thetable{} continued.}\\
\toprule
 &  &  &  & \multicolumn{3}{c}{Before matching} & \multicolumn{3}{c}{After matching} \\
\cmidrule(lr){5-7}\cmidrule(lr){8-10}
Platform & Design & Snapshot & Covariate & T mean & C mean & SMD & T mean & C mean & SMD \\
\midrule
\endhead
\bottomrule
\endfoot
Spotify & Hist. & 2016 & Chart debut year & 1983.17 & 1990.20 & -0.391 & 1983.17 & 1983.19 & -0.001 \\
Spotify & Hist. & 2016 & Song age & 33.63 & 26.60 & 0.391 & 33.63 & 33.61 & 0.001 \\
Spotify & Hist. & 2016 & Billboard weeks & 11.20 & 12.06 & -0.108 & 11.20 & 11.14 & 0.009 \\
Spotify & Hist. & 2016 & Peak chart position & 44.97 & 47.38 & -0.081 & 44.97 & 44.65 & 0.011 \\
Spotify & Hist. & 2016 & Artist Billboard catalog size & 20.49 & 10.01 & 0.451 & 20.49 & 18.55 & 0.121 \\
Spotify & Hist. & 2016 & Artist total Billboard weeks, excluding focal song & 222.51 & 60.02 & 1.074 & 222.51 & 201.09 & 0.125 \\
Spotify & Hist. & 2016 & Artist mean Billboard weeks, excluding focal song & 12.20 & 10.64 & 0.287 & 12.20 & 12.18 & 0.006 \\
Spotify & Hist. & 2016 & Artist best other-song peak position & 6.95 & 26.02 & -0.922 & 6.95 & 7.32 & -0.031 \\
Spotify & Hist. & 2016 & Artist superstar indicator, top 1 percent & 0.23 & 0.05 & 0.526 & 0.23 & 0.23 & 0.001 \\
Spotify & Hist. & 2017 & Chart debut year & 1982.57 & 1986.55 & -0.223 & 1982.57 & 1982.68 & -0.006 \\
Spotify & Hist. & 2017 & Song age & 34.98 & 31.00 & 0.223 & 34.98 & 34.87 & 0.006 \\
Spotify & Hist. & 2017 & Billboard weeks & 10.74 & 10.89 & -0.021 & 10.74 & 10.67 & 0.009 \\
Spotify & Hist. & 2017 & Peak chart position & 46.31 & 51.03 & -0.161 & 46.31 & 46.07 & 0.009 \\
Spotify & Hist. & 2017 & Artist Billboard catalog size & 21.65 & 11.68 & 0.364 & 21.65 & 18.94 & 0.129 \\
Spotify & Hist. & 2017 & Artist total Billboard weeks, excluding focal song & 225.67 & 69.44 & 0.938 & 225.67 & 196.42 & 0.159 \\
Spotify & Hist. & 2017 & Artist mean Billboard weeks, excluding focal song & 11.99 & 10.51 & 0.275 & 11.99 & 11.99 & -0.000 \\
Spotify & Hist. & 2017 & Artist best other-song peak position & 6.83 & 27.24 & -0.974 & 6.83 & 7.27 & -0.039 \\
Spotify & Hist. & 2017 & Artist superstar indicator, top 1 percent & 0.22 & 0.05 & 0.504 & 0.22 & 0.22 & 0.012 \\
Spotify & Prior & 2017 & Chart debut year & 1983.24 & 1988.04 & -0.272 & 1983.24 & 1983.55 & -0.018 \\
Spotify & Prior & 2017 & Song age & 34.31 & 29.51 & 0.272 & 34.31 & 34.00 & 0.018 \\
Spotify & Prior & 2017 & Billboard weeks & 11.22 & 11.43 & -0.027 & 11.22 & 11.29 & -0.011 \\
Spotify & Prior & 2017 & Peak chart position & 44.94 & 49.52 & -0.156 & 44.94 & 44.35 & 0.020 \\
Spotify & Prior & 2017 & Artist Billboard catalog size & 20.47 & 12.43 & 0.315 & 20.47 & 17.84 & 0.163 \\
Spotify & Prior & 2017 & Artist total Billboard weeks, excluding focal song & 222.68 & 74.99 & 0.959 & 222.68 & 194.74 & 0.162 \\
Spotify & Prior & 2017 & Artist mean Billboard weeks, excluding focal song & 12.22 & 10.75 & 0.272 & 12.22 & 12.19 & 0.008 \\
Spotify & Prior & 2017 & Artist best other-song peak position & 6.96 & 26.40 & -0.946 & 6.96 & 7.92 & -0.082 \\
Spotify & Prior & 2017 & Artist superstar indicator, top 1 percent & 0.23 & 0.06 & 0.484 & 0.23 & 0.23 & 0.001 \\
Spotify & Prior & 2017 & Prior Spotify popularity & 29.48 & 29.40 & 0.004 & 29.48 & 29.29 & 0.010 \\
Spotify & Hist. & 2022 & Chart debut year & 1985.13 & 1989.26 & -0.209 & 1985.13 & 1985.35 & -0.011 \\
Spotify & Hist. & 2022 & Song age & 37.52 & 33.40 & 0.209 & 37.52 & 37.30 & 0.011 \\
Spotify & Hist. & 2022 & Billboard weeks & 10.40 & 10.66 & -0.034 & 10.40 & 10.38 & 0.003 \\
Spotify & Hist. & 2022 & Peak chart position & 46.98 & 51.01 & -0.138 & 46.98 & 47.07 & -0.003 \\
Spotify & Hist. & 2022 & Artist Billboard catalog size & 24.77 & 12.20 & 0.421 & 24.77 & 20.78 & 0.154 \\
Spotify & Hist. & 2022 & Artist total Billboard weeks, excluding focal song & 238.94 & 79.89 & 0.825 & 238.94 & 216.22 & 0.104 \\
Spotify & Hist. & 2022 & Artist mean Billboard weeks, excluding focal song & 11.77 & 10.30 & 0.273 & 11.77 & 11.90 & -0.026 \\
Spotify & Hist. & 2022 & Artist best other-song peak position & 6.86 & 25.06 & -0.879 & 6.86 & 7.01 & -0.013 \\
Spotify & Hist. & 2022 & Artist superstar indicator, top 1 percent & 0.23 & 0.07 & 0.459 & 0.23 & 0.23 & 0.000 \\
Spotify & Prior & 2022 & Chart debut year & 1982.85 & 1986.97 & -0.231 & 1982.85 & 1983.16 & -0.018 \\
Spotify & Prior & 2022 & Song age & 39.80 & 35.68 & 0.231 & 39.80 & 39.49 & 0.018 \\
Spotify & Prior & 2022 & Billboard weeks & 10.78 & 11.02 & -0.032 & 10.78 & 10.86 & -0.011 \\
Spotify & Prior & 2022 & Peak chart position & 46.17 & 50.83 & -0.159 & 46.17 & 45.67 & 0.018 \\
Spotify & Prior & 2022 & Artist Billboard catalog size & 21.78 & 11.88 & 0.357 & 21.78 & 18.50 & 0.151 \\
Spotify & Prior & 2022 & Artist total Billboard weeks, excluding focal song & 226.88 & 71.13 & 0.925 & 226.88 & 189.63 & 0.201 \\
Spotify & Prior & 2022 & Artist mean Billboard weeks, excluding focal song & 12.03 & 10.60 & 0.264 & 12.03 & 11.97 & 0.014 \\
Spotify & Prior & 2022 & Artist best other-song peak position & 6.85 & 27.13 & -0.969 & 6.85 & 7.69 & -0.073 \\
Spotify & Prior & 2022 & Artist superstar indicator, top 1 percent & 0.23 & 0.06 & 0.500 & 0.23 & 0.22 & 0.018 \\
Spotify & Prior & 2022 & Prior Spotify popularity & 32.92 & 30.88 & 0.102 & 32.92 & 32.75 & 0.009 \\
Spotify & Hist. & 2025 & Chart debut year & 1986.31 & 1992.54 & -0.291 & 1986.31 & 1986.60 & -0.014 \\
Spotify & Hist. & 2025 & Song age & 39.34 & 33.11 & 0.291 & 39.34 & 39.05 & 0.014 \\
Spotify & Hist. & 2025 & Billboard weeks & 10.12 & 10.12 & 0.000 & 10.12 & 10.09 & 0.004 \\
Spotify & Hist. & 2025 & Peak chart position & 47.05 & 51.80 & -0.163 & 47.05 & 46.96 & 0.003 \\
Spotify & Hist. & 2025 & Artist Billboard catalog size & 27.23 & 12.42 & 0.458 & 27.23 & 21.96 & 0.178 \\
Spotify & Hist. & 2025 & Artist total Billboard weeks, excluding focal song & 252.66 & 83.58 & 0.788 & 252.66 & 226.04 & 0.111 \\
Spotify & Hist. & 2025 & Artist mean Billboard weeks, excluding focal song & 11.60 & 10.15 & 0.264 & 11.60 & 11.77 & -0.035 \\
Spotify & Hist. & 2025 & Artist best other-song peak position & 6.69 & 24.40 & -0.867 & 6.69 & 6.80 & -0.010 \\
Spotify & Hist. & 2025 & Artist superstar indicator, top 1 percent & 0.24 & 0.08 & 0.474 & 0.24 & 0.24 & 0.000 \\
Spotify & Prior & 2025 & Chart debut year & 1985.19 & 1989.25 & -0.205 & 1985.19 & 1985.41 & -0.011 \\
Spotify & Prior & 2025 & Song age & 40.46 & 36.40 & 0.205 & 40.46 & 40.24 & 0.011 \\
Spotify & Prior & 2025 & Billboard weeks & 10.42 & 10.72 & -0.037 & 10.42 & 10.36 & 0.007 \\
Spotify & Prior & 2025 & Peak chart position & 46.87 & 50.84 & -0.136 & 46.87 & 46.93 & -0.002 \\
Spotify & Prior & 2025 & Artist Billboard catalog size & 24.79 & 12.68 & 0.401 & 24.79 & 20.42 & 0.168 \\
Spotify & Prior & 2025 & Artist total Billboard weeks, excluding focal song & 239.35 & 83.16 & 0.806 & 239.35 & 210.66 & 0.132 \\
Spotify & Prior & 2025 & Artist mean Billboard weeks, excluding focal song & 11.79 & 10.38 & 0.262 & 11.79 & 11.85 & -0.013 \\
Spotify & Prior & 2025 & Artist best other-song peak position & 6.84 & 24.61 & -0.863 & 6.84 & 7.17 & -0.029 \\
Spotify & Prior & 2025 & Artist superstar indicator, top 1 percent & 0.23 & 0.08 & 0.447 & 0.23 & 0.23 & 0.003 \\
Spotify & Prior & 2025 & Prior Spotify popularity & 37.73 & 37.08 & 0.030 & 37.73 & 37.69 & 0.002 \\
Last.fm & Prior & 2017 & Chart debut year & 1983.22 & 1988.02 & -0.272 & 1983.22 & 1983.52 & -0.018 \\
Last.fm & Prior & 2017 & Song age & 34.33 & 29.53 & 0.272 & 34.33 & 34.03 & 0.018 \\
Last.fm & Prior & 2017 & Billboard weeks & 11.21 & 11.42 & -0.029 & 11.21 & 11.29 & -0.012 \\
Last.fm & Prior & 2017 & Peak chart position & 44.94 & 49.55 & -0.158 & 44.94 & 44.35 & 0.020 \\
Last.fm & Prior & 2017 & Artist Billboard catalog size & 20.50 & 12.52 & 0.310 & 20.50 & 17.86 & 0.165 \\
Last.fm & Prior & 2017 & Artist total Billboard weeks, excluding focal song & 222.90 & 75.23 & 0.958 & 222.90 & 194.95 & 0.162 \\
Last.fm & Prior & 2017 & Artist mean Billboard weeks, excluding focal song & 12.22 & 10.74 & 0.270 & 12.22 & 12.18 & 0.007 \\
Last.fm & Prior & 2017 & Artist best other-song peak position & 6.93 & 26.47 & -0.951 & 6.93 & 7.89 & -0.082 \\
Last.fm & Prior & 2017 & Artist superstar indicator, top 1 percent & 0.23 & 0.06 & 0.484 & 0.23 & 0.23 & 0.001 \\
Last.fm & Prior & 2017 & Prior Spotify popularity & 29.49 & 29.38 & 0.006 & 29.49 & 29.31 & 0.010 \\
\end{longtable}
\begin{flushleft}\footnotesize Hist.\ = historical matching; Prior = prior-attention matching. Before columns compare all catalog-exposed with all unexposed non-film songs before matching; after columns compare matched samples. Artist catalogs vary strongly in historical visibility and size, which is why the full covariate-level balance is retained. \end{flushleft}
}


{\footnotesize
\setlength{\tabcolsep}{2.5pt}
\renewcommand{\arraystretch}{1.05}
\begin{longtable}{>{\raggedright\arraybackslash}p{0.19\linewidth}>{\raggedright\arraybackslash}p{0.17\linewidth}>{\centering\arraybackslash}p{0.052\linewidth}>{\centering\arraybackslash}p{0.078\linewidth}>{\centering\arraybackslash}p{0.082\linewidth}>{\centering\arraybackslash}p{0.052\linewidth}>{\centering\arraybackslash}p{0.082\linewidth}>{\centering\arraybackslash}p{0.105\linewidth}}
\caption{Replacement-aware inference and design sensitivity for the matched catalog premium.}\label{tab:si-catalog-matching-inference}\\
\toprule
Comparison & Variant & Pairs & \makecell[b]{Control\\ IDs} & Estimate & \makecell[b]{Naive\\ SE} & \makecell[b]{Cluster\\ SE\\ (control)} & 95\% CI \\
\midrule
\endfirsthead
\multicolumn{8}{@{}l}{\footnotesize Table~\thetable{} continued.}\\
\toprule
Comparison & Variant & Pairs & \makecell[b]{Control\\ IDs} & Estimate & \makecell[b]{Naive\\ SE} & \makecell[b]{Cluster\\ SE\\ (control)} & 95\% CI \\
\midrule
\endhead
\bottomrule
\endfoot
catalog baseline-attention | 2017-07 (with replacement) & with replacement & 6,794 & 1,927 & 0.46 & 0.087 & 0.174 & [0.11, 0.79] \\
catalog baseline-attention | 2017-07 (no replacement) & no replacement & 3,271 & 3,271 & 1.29 & 0.130 & -- & [1.04, 1.55] \\
catalog baseline-attention | 2017-07 (stricter caliper 2.0) & stricter caliper & 6,794 & 2,019 & 0.26 & 0.081 & -- & [0.10, 0.42] \\
catalog baseline-attention | 2017-07 (CEM decade x baseline-pop quintile) & CEM & 6,794 & -- & 1.57 & -- & -- & -- \\
catalog baseline-attention | 2022-08 (with replacement) & with replacement & 7,712 & 2,197 & -0.31 & 0.131 & 0.317 & [-0.97, 0.26] \\
catalog baseline-attention | 2022-08 (no replacement) & no replacement & 3,740 & 3,740 & 0.32 & 0.183 & -- & [-0.04, 0.67] \\
catalog baseline-attention | 2022-08 (stricter caliper 2.0) & stricter caliper & 7,712 & 2,293 & -0.35 & 0.127 & -- & [-0.60, -0.10] \\
catalog baseline-attention | 2022-08 (CEM decade x baseline-pop quintile) & CEM & 7,712 & -- & -0.70 & -- & -- & -- \\
catalog baseline-attention | 2025 (with replacement) & with replacement & 9,285 & 3,049 & 0.44 & 0.091 & 0.184 & [0.09, 0.81] \\
catalog baseline-attention | 2025 (no replacement) & no replacement & 5,268 & 5,265 & 1.02 & 0.124 & -- & [0.78, 1.27] \\
catalog baseline-attention | 2025 (stricter caliper 2.0) & stricter caliper & 9,285 & 3,186 & 0.48 & 0.087 & -- & [0.31, 0.65] \\
catalog baseline-attention | 2025 (CEM decade x baseline-pop quintile) & CEM & 9,285 & -- & 1.17 & -- & -- & -- \\
\end{longtable}
\begin{flushleft}\footnotesize Cluster-by-control intervals account for control reuse in the prior-attention catalog comparison. Under this inference, the August 2022 contrast is not distinguishable from zero. Its sign varies across matching designs: the no-replacement estimate is positive, whereas the stricter-caliper and coarsened-exact-matching estimates are negative. Control IDs are alphanumeric song-artist identifiers. No-replacement matching excludes repeated record indices, not different records sharing an ID. \end{flushleft}
}


Table~\ref{tab:si-catalog-retention} reports matched catalog attention changes. The prior-attention matched pairs are held fixed, retaining the pair sets, covariates, five-point prior-attention caliper, and cohort support of Table~\ref{tab:si-catalog-matching}. Each pair is re-expressed as prior-to-current attention changes, and the difference in change between exposed songs and their matched controls is the estimand. No rematching is performed. Because catalog exposure is defined at the artist level, matched pairs are dependent within artist identities on both sides. The primary interval is a two-way pigeonhole bootstrap that independently resamples treated artist identities and reused control artist identities and reweights pairs by the product of the artist-cluster draw counts (2,000 deterministic-seed replications). Cohort-specific estimates vary in sign, and every artist-level interval includes zero. The pooled descriptive summary is small and its interval includes zero. Leave-one-cohort-out combinations (Table~\ref{tab:si-catalog-retention-loco}) remain inconclusive in every configuration. These comparisons do not establish a corresponding attention increase in other Billboard songs by the same artists. They also do not establish the absence of such an increase or of attention spreading beyond the observed songs.

\begin{table}[!htbp]
\centering
\footnotesize
\caption{Catalog attention changes on the fixed prior-attention matched pairs.}
\label{tab:si-catalog-retention}
\renewcommand{\arraystretch}{1.1}
\setlength{\tabcolsep}{3pt}
\begin{tabular}{>{\raggedright\arraybackslash}p{0.105\linewidth}>{\centering\arraybackslash}p{0.062\linewidth}>{\centering\arraybackslash}p{0.072\linewidth}>{\centering\arraybackslash}p{0.072\linewidth}>{\centering\arraybackslash}p{0.072\linewidth}>{\centering\arraybackslash}p{0.072\linewidth}>{\centering\arraybackslash}p{0.072\linewidth}>{\centering\arraybackslash}p{0.105\linewidth}>{\raggedright\arraybackslash}p{0.18\linewidth}}
\toprule
Cohort & Pairs & \makecell[b]{Treated\\ artists} & \makecell[b]{Control\\ artists} & \makecell[b]{Treated\\ change} & \makecell[b]{Control\\ change} & \makecell[b]{Diff.\ in\\ change} & \makecell[b]{Primary\\ 95\% CI} & Observed change pattern \\
\midrule
July 2017 & 6,794 & 1,293 & 802 & 4.63 & 4.36 & 0.27 & [-0.36, 0.89] & treated and controls both increased; treated increased slightly more \\
August 2022 & 7,712 & 1,358 & 889 & 3.41 & 3.89 & -0.48 & [-1.78, 0.71] & treated increased less than controls \\
August 2025 & 9,285 & 1,485 & 1,053 & 0.03 & -0.36 & 0.40 & [-0.22, 0.95] & treated remained approximately stable; controls declined \\
Pooled (fixed-effect, descriptive) & 23,791 & -- & -- & -- & -- & 0.25 & [-0.15, 0.66] & descriptive pooled interval includes zero \\
\bottomrule
\end{tabular}
\begin{flushleft}\footnotesize
Prior-to-current change scores on the identical prior-attention matched catalog pairs; no rematching is performed. The primary interval is a two-way pigeonhole bootstrap over treated and reused control artist identities (2{,}000 deterministic-seed replications), because catalog exposure is defined at the artist level and pairs are dependent within artists on both sides. The pooled row is a descriptive inverse-variance combination of three heterogeneous cohorts and should not override the cohort-specific pattern. Pattern descriptions summarize observed group changes and are descriptive; the artist-level confidence intervals provide the inferential assessment.
\end{flushleft}
\end{table}


\begin{table}[!htbp]
\centering
\footnotesize
\caption{Leave-one-cohort-out sensitivity of the pooled catalog attention-change summary.}
\label{tab:si-catalog-retention-loco}
\renewcommand{\arraystretch}{1.05}
\begin{tabular}{llll}
\toprule
Omitted cohort & Included cohorts & Pooled diff.\ in change & 95\% CI \\
\midrule
July 2017 & August 2022, August 2025 & 0.24 & [-0.30, 0.77] \\
August 2022 & July 2017, August 2025 & 0.34 & [-0.09, 0.76] \\
August 2025 & July 2017, August 2022 & 0.12 & [-0.43, 0.68] \\
\bottomrule
\end{tabular}
\begin{flushleft}\footnotesize
Each row recomputes the descriptive fixed-effect pooled difference-in-change after omitting one cohort, using the primary artist-level variances. With only three heterogeneous cohorts these are sensitivity descriptions, not formal heterogeneity tests.
\end{flushleft}
\end{table}


\section*{Inference and Robustness Checks}

Pooled models reuse the same Billboard song-artist record across snapshots. Table~\ref{tab:si-cluster-inference} refits every pooled specification exactly as reported in the paper (time-respecting exposures where those are primary) and compares HC1 with cluster-by-song standard errors; no pooled conclusion changes under cluster-by-song inference. Table~\ref{tab:si-robustness-summary} consolidates the robustness and placebo checks: the future-placement placebo (selection before observed placement), the superstar exclusion, and the collaboration-string exclusion for the catalog analysis.

{\footnotesize
\setlength{\tabcolsep}{2.5pt}
\renewcommand{\arraystretch}{1.05}
\begin{longtable}{>{\raggedright\arraybackslash}p{0.19\linewidth}>{\raggedright\arraybackslash}p{0.15\linewidth}>{\centering\arraybackslash}p{0.055\linewidth}>{\centering\arraybackslash}p{0.062\linewidth}>{\centering\arraybackslash}p{0.082\linewidth}>{\centering\arraybackslash}p{0.045\linewidth}>{\centering\arraybackslash}p{0.045\linewidth}>{\centering\arraybackslash}p{0.1\linewidth}>{\centering\arraybackslash}p{0.072\linewidth}}
\caption{Cluster-by-song inference for pooled models with repeated song-snapshot observations.}\label{tab:si-cluster-inference}\\
\toprule
 &  & \multicolumn{2}{c}{Support} &  & \multicolumn{1}{c}{HC1} & \multicolumn{3}{c}{Cluster-by-song inference} \\
\cmidrule(lr){3-4}\cmidrule(lr){7-9}
Model & Term & N obs. & N clusters & Estimate & SE & SE & 95\% CI & $p$ \\
\midrule
\endfirsthead
\multicolumn{9}{@{}l}{\footnotesize Table~\thetable{} continued.}\\
\toprule
 &  & \multicolumn{2}{c}{Support} &  & \multicolumn{1}{c}{HC1} & \multicolumn{3}{c}{Cluster-by-song inference} \\
\cmidrule(lr){3-4}\cmidrule(lr){7-9}
Model & Term & N obs. & N clusters & Estimate & SE & SE & 95\% CI & $p$ \\
\midrule
\endhead
\bottomrule
\endfoot
Pooled broad Spotify film-linked premium & film linked & 83,785 & 28,281 & 13.79 & 0.128 & 0.227 & [13.35, 14.24] & $<0.001$ \\
Pooled standardized Spotify film-linked premium & film linked & 83,785 & 28,281 & 0.66 & 0.006 & 0.011 & [0.63, 0.68] & $<0.001$ \\
Repeated embedding (time-respecting, primary) & 1 realized films vs none & 83,780 & 28,281 & 9.25 & 0.153 & 0.264 & [8.73, 9.76] & $<0.001$ \\
Repeated embedding (time-respecting, primary) & 2-3 realized films vs none & 83,780 & 28,281 & 16.72 & 0.203 & 0.341 & [16.05, 17.39] & $<0.001$ \\
Repeated embedding (time-respecting, primary) & 4+ realized films vs none & 83,780 & 28,281 & 24.19 & 0.249 & 0.424 & [23.35, 25.02] & $<0.001$ \\
Film visibility (continuous, time-respecting, primary) & log1p max recorded worldwide box office; films released by snapshot year & 15,245 & 4,294 & 1.15 & 0.058 & 0.102 & [0.95, 1.35] & $<0.001$ \\
Song event-time model (broad) & pre (<= -2 years) & 83,785 & 28,281 & 9.50 & 0.671 & 0.865 & [7.80, 11.19] & $<0.001$ \\
Song event-time model (broad) & event year & 83,785 & 28,281 & 11.23 & 0.908 & 0.911 & [9.44, 13.02] & $<0.001$ \\
Song event-time model (broad) & 1-5 years post & 83,785 & 28,281 & 10.16 & 0.357 & 0.436 & [9.31, 11.02] & $<0.001$ \\
Song event-time model (broad) & >5 years post & 83,785 & 28,281 & 14.18 & 0.136 & 0.239 & [13.71, 14.65] & $<0.001$ \\
Pooled artist-catalog regression (final specification) & artist has treated song & 57,445 & 19,067 & 2.47 & 0.126 & 0.209 & [2.06, 2.88] & $<0.001$ \\
\end{longtable}
\begin{flushleft}\footnotesize For the pooled models listed here, the table compares HC1 and cluster-by-song standard errors using the same specification and exposure definition. Clusters are defined by alphanumeric song-title and artist identifiers, not by exact Billboard record strings. Multiple exact records can therefore share a cluster. The listed estimates retain their sign and exclude zero under cluster-by-song inference. These comparisons address repeated song-snapshot dependence and do not change the estimands. \end{flushleft}
}


\begin{table}[!htbp]
\centering
\footnotesize
\caption{Robustness and placebo checks (exact final specifications, cluster-by-song inference).}
\label{tab:si-robustness-summary}
\renewcommand{\arraystretch}{1.1}
\setlength{\tabcolsep}{2.5pt}
\begin{tabular}{>{\raggedright\arraybackslash}p{0.13\linewidth}>{\raggedright\arraybackslash}p{0.15\linewidth}>{\centering\arraybackslash}p{0.082\linewidth}>{\centering\arraybackslash}p{0.07\linewidth}>{\centering\arraybackslash}p{0.085\linewidth}>{\centering\arraybackslash}p{0.075\linewidth}>{\centering\arraybackslash}p{0.065\linewidth}>{\centering\arraybackslash}p{0.078\linewidth}>{\raggedright\arraybackslash}p{0.15\linewidth}}
\toprule
Check & Exact specification & Estimate & \makecell[b]{Cluster\\ SE} & 95\% CI & $p$ & N obs. & \makecell[b]{N\\ clusters} & Interpretation \\
\midrule
Future-placement placebo & film-linked song observed before its first film appearance & 9.57 & 0.665 & [8.27, 10.88] & $<0.001$ & 67,453 & 23,994 & attention difference before observed film placement \\
Superstar exclusion (pooled broad premium) & ever film-linked; songs with >= 40 Billboard weeks excluded (top 1\%) & 13.95 & 0.230 & [13.50, 14.40] & $<0.001$ & 83,002 & 27,999 & positive association after excluding superstar songs \\
Catalog premium excluding collaboration-like artist strings & artist has another film-linked Billboard song; collaboration-like artist strings excluded & 3.21 & 0.223 & [2.78, 3.65] & $<0.001$ & 52,967 & 17,155 & positive association after excluding collaboration-like strings \\
\bottomrule
\end{tabular}
\begin{flushleft}\footnotesize
The future-placement placebo compares not-yet-film-linked song-snapshots with never-film songs; a positive value documents an attention difference before observed placement. The superstar exclusion refits the pooled broad premium without the top 1\% of songs by Billboard weeks. The collaboration-string exclusion refits the exact final pooled catalog specification after removing collaboration-like artist strings, addressing catalog-key sensitivity.
\end{flushleft}
\end{table}


\end{document}
